% =====================================================================
% v2 (event-episode) experiment: methodology + findings.
% Intended to be \input{}-ed into the main progress_report.tex once
% integrated.  Requires:
%   \usepackage{booktabs,longtable,graphicx,amsmath,amssymb}
% All numerical tables are auto-generated by
%   scripts/v2/make_v2_latex_tables.py
% from the CSV artefacts in v2_results_second_run/v2/tables/.
% =====================================================================

\section{Event-Episode Re-Evaluation (v2)}
\label{sec:v2}

The original analysis (Sections~\ref{sec:exp1}--\ref{sec:exp2}) selected
hyperparameters on the held-out test split, fused training and testing
labels under a single weekly-OR-of-flags scheme, and reported metrics
across five random patient seeds without an isolated validation set.
Reviewer feedback identified three concrete weaknesses: (i)~the
labelling target was clinically heterogeneous (any worsening flag
counted as a positive week), (ii)~hyperparameter optimisation and
evaluation shared the same partition, and (iii)~there was no
prospective check that the apparently positive signal would survive a
genuinely unseen patient cohort.

This section addresses all three with a frozen, non-destructive
re-implementation that we refer to as \emph{v2}.  No file from the
v1 pipeline was modified; all new code lives under
\texttt{src/event\_v2/}, all new scripts under \texttt{scripts/v2/},
and all artefacts under \texttt{outputs/v2/tables/}.  The v2 task
operates on a more clinically meaningful binary target and reports
metrics on a patient-disjoint test set that was never touched during
hyperparameter selection.  Every per-cell, per-algorithm, per-subset
and per-seed number reported below is reproduced verbatim from those
CSVs through \texttt{scripts/v2/make\_v2\_latex\_tables.py}; nothing
in this section is hand-typed.

\subsection{Task definition}
\label{sec:v2-task}

For every event-episode~$E$ (a calendar day on which the
threshold~$T$ for the worsening score is exceeded, with~$T\!\in\!\{2,3,4\}$
representing increasingly severe events), we define a binary
classification problem on per-day samples:
\[
y(d) =
\begin{cases}
1, & d \in [E-L,\, E-1] \text{ for some event } E, \\
0, & d \text{ lies in a stable week with no event, no positive-week} \\[-1pt]
   & \quad\text{overlap, and no day-level washout overlap with any event.}
\end{cases}
\]
The lookback length~$L\!\in\!\{3,7,14\}$ days controls how far before the
event positives are emitted; the washout~$W\!\in\!\{0,7,14\}$ days
controls how many days around an event are excluded from the negative
pool.  By construction, the event day~$E$ itself is never a positive
sample, so the model is asked to predict a forthcoming episode rather
than to detect one on the day it occurs.  The grid
$(T,L,W)\!\in\!\{2,3,4\}\!\times\!\{3,7,14\}\!\times\!\{0,7,14\}$
gives 27~task instances, all evaluated end-to-end.

The resulting prevalence is severely imbalanced (between $0.6\%$
and $4.5\%$ positives depending on the cell), which is by design
closer to the operational scenario of predicting a worsening episode
than the $30.3\%$ weekly-OR prevalence used in v1.  Per-cell sample
counts and exclusion reasons are listed in
Table~\ref{tab:v2_sample_counts}.

\begin{table}[t]
\centering
\small
\caption{Sample counts for every $(T,L,W)$ cell across the entire cohort.  Excluded counts explain how the positive and negative pools are pruned.}\label{tab:v2_sample_counts}
\begin{tabular}{ccccccccc}
\toprule
$T$ & $L$ & $W$ & Pos cand & Pos kept & Pos exc (hist) & Neg cand & Neg kept & Neg exc (total) \\
\midrule
2 & 3 & 0 & 63 & 59 & 4 & 1580 & 1529 & 51 \\
2 & 3 & 7 & 63 & 59 & 4 & 1580 & 1419 & 161 \\
2 & 3 & 14 & 63 & 59 & 4 & 1580 & 1323 & 257 \\
2 & 7 & 0 & 63 & 56 & 7 & 1517 & 1388 & 129 \\
2 & 7 & 7 & 63 & 56 & 7 & 1517 & 1290 & 227 \\
2 & 7 & 14 & 63 & 56 & 7 & 1517 & 1221 & 296 \\
2 & 14 & 0 & 63 & 54 & 9 & 1419 & 1191 & 228 \\
2 & 14 & 7 & 63 & 54 & 9 & 1419 & 1130 & 289 \\
2 & 14 & 14 & 63 & 54 & 9 & 1419 & 1067 & 352 \\
3 & 3 & 0 & 63 & 59 & 4 & 1580 & 1529 & 51 \\
3 & 3 & 7 & 63 & 59 & 4 & 1580 & 1425 & 155 \\
3 & 3 & 14 & 63 & 59 & 4 & 1580 & 1327 & 253 \\
3 & 7 & 0 & 63 & 55 & 8 & 1517 & 1388 & 129 \\
3 & 7 & 7 & 63 & 55 & 8 & 1517 & 1297 & 220 \\
3 & 7 & 14 & 63 & 55 & 8 & 1517 & 1225 & 292 \\
3 & 14 & 0 & 63 & 53 & 10 & 1419 & 1191 & 228 \\
3 & 14 & 7 & 63 & 53 & 10 & 1419 & 1134 & 285 \\
3 & 14 & 14 & 63 & 53 & 10 & 1419 & 1071 & 348 \\
4 & 3 & 0 & 63 & 60 & 3 & 1580 & 1529 & 51 \\
4 & 3 & 7 & 63 & 60 & 3 & 1580 & 1430 & 150 \\
4 & 3 & 14 & 63 & 60 & 3 & 1580 & 1331 & 249 \\
4 & 7 & 0 & 63 & 57 & 6 & 1517 & 1388 & 129 \\
4 & 7 & 7 & 63 & 57 & 6 & 1517 & 1302 & 215 \\
4 & 7 & 14 & 63 & 57 & 6 & 1517 & 1230 & 287 \\
4 & 14 & 0 & 63 & 55 & 8 & 1419 & 1191 & 228 \\
4 & 14 & 7 & 63 & 55 & 8 & 1419 & 1139 & 280 \\
4 & 14 & 14 & 63 & 55 & 8 & 1419 & 1076 & 343 \\
\bottomrule
\end{tabular}
\end{table}


\subsection{Three-way patient split}
\label{sec:v2-split}

We replaced the v1 two-way patient split with a deterministic
three-way patient-disjoint split using
$\texttt{seed}\!=\!42$:
\[
\text{train}:\text{val}:\text{test} = 70:15:20.
\]
With the AAMOS-00 cohort of $18$~patients, this assigns
$13$~patients to training, $3$~to validation, and $4$~to test.
The split is computed once and persisted to
\texttt{outputs/v2/tables/tune\_tabular\_v2\_split.json}; every
subsequent phase consumes the same file so that no information about
test-patient identities leaks into hyperparameter selection.

\subsection{Hyperparameter optimisation}
\label{sec:v2-hpo}

\paragraph{Tabular HPO (Phase~4).}
For each $(T,L,W)$ cell we sweep three families:
logistic regression ($C\!\in\!\{10^{-2},10^{-1},1,10\}\times
\{\text{none}, \text{balanced}\}$, $8$~combinations);
random forest ($n_{\text{est}}\!\in\!\{200,500\}\times
\max\_\text{depth}\!\in\!\{6,10,\infty\}\times
\min\_\text{leaf}\!\in\!\{1,5\}\times\{\text{balanced}\}$,
$24$~combinations);
XGBoost ($n_{\text{est}}\!\in\!\{200,500\}\times
\max\_\text{depth}\!\in\!\{3,6\}\times
\text{lr}\!\in\!\{0.05,0.1\}\times
\text{subsample}\!\in\!\{0.7,1.0\}\times
\text{scale\_pos\_weight}\!\in\!\{1,\text{auto}\}$,
$32$~combinations).
All models are fit on the training split and scored on the validation
split; the winner per cell per algorithm is the one with highest
validation PR-AUC.  The full trial table is in
\texttt{tune\_tabular\_v2\_trials.csv} and the winners in
\texttt{tune\_tabular\_v2\_best.csv}.

\paragraph{Deep-learning HPO (Phase~5).}
The same $(T,L,W)$ grid is replayed for four sequence architectures:
GRU, LSTM, RNN, CNN.  Inputs are daily-level $L\!\times\!D$ tensors
with $D\!=\!72$ feature channels formed from all available sensor
sources (smartwatch, smart inhaler, peak flow, environment).
The coarse architecture grid is hidden dimension
$\in\!\{32,64\}$, layers $\in\!\{1,2\}$, dropout
$\in\!\{0.2,0.4\}$ for recurrent models and dropout
$\in\!\{0.2,0.4\}$ for the CNN.  The overall PR-AUC winner is
re-fit across five seeds ($42,43,44,45,46$) to estimate seed
sensitivity; outputs are written to
\texttt{tune\_dl\_v2\_multiseed.csv}.

Throughout, only the training and validation partitions are exposed
to the HPO loop.  The test partition is sealed until Phase~8.

\subsection{Sensor ablation}
\label{sec:v2-ablation}

For each $(T,L,W)$ cell with a successful tabular winner, we refit the
winner's hyperparameters on nine sensor subsets: the full
(\texttt{all}) channel pool, four singletons
(\texttt{smartwatch\_only}, \texttt{smartinhaler\_only},
\texttt{peakflow\_only}, \texttt{environment\_only}), and four
leave-one-out variants (\texttt{no\_smartwatch}, etc.).  Both the
validation and test metrics are recorded so we can ask whether
multi-modal fusion genuinely outperforms any single source.  Results
are in \texttt{sensor\_ablation\_v2.csv} ($27\times 3\times 9 = 729$
rows); the complete grid is reproduced in
Table~\ref{tab:v2_sensor_full} and the top-3 subsets per algorithm
in Table~\ref{tab:v2_ablation_topk}.

\subsection{Leakage probe}
\label{sec:v2-leakage}

Before any test-set prediction is made we run a deliberate shuffle
test.  For each $(T,L,W,\text{algo})$ winner, the training labels
are permuted under five independent random number generators while
features and validation labels are kept intact; the model is then
re-fit and its validation ROC-AUC and PR-AUC are recorded
(\texttt{leakage\_probe\_v2.csv}, summary in
\texttt{leakage\_probe\_v2\_summary.csv}; full table in
Table~\ref{tab:v2_leakage}).  If the pipeline is leakage-free, the
\emph{average} shuffled validation ROC-AUC across all configurations
should sit near $0.5$.  In our run the overall mean is $0.505$ with
standard deviation $0.125$ across the $81$ shuffled-label
configurations; the offline analysis script reports
\texttt{OK: overall mean is near 0.5 (no systematic leakage)}.
Per-cell deviations from $0.5$ are expected with $n_{\text{val}}^+
\!\in\!\{1,3\}$ positives and are not actionable in isolation.

\subsection{Final test evaluation}
\label{sec:v2-test}

In Phase~8 each tabular winner is refit on the union of training and
validation data with the hyperparameters chosen in Phase~4 and
applied once to the test partition.  Per-sample predictions are
saved to \texttt{final\_test\_v2\_predictions\_tabular.csv} (one row
per test sample per winning configuration).  The deep-learning
winner is refit across the five seeds used in Phase~5; per-sample,
per-seed predictions land in
\texttt{final\_test\_v2\_predictions\_dl.csv}.

Because the test class prevalence is very low we additionally
estimate a probability cutoff by maximising F1 on the validation
split using a train-only fit (to avoid using the val partition
twice).  The tuned threshold~$\tau^*$ is reported alongside the
default $0.5$ cutoff in \texttt{final\_test\_v2\_summary.csv} and in
all results tables below.

\subsection{Bootstrap analysis}
\label{sec:v2-bootstrap}

To quantify the considerable uncertainty introduced by the small
number of test positives, we compute $2{,}000$-resample bootstrap
$95\%$ confidence intervals on test PR-AUC and ROC-AUC for every
tabular winner and for the DL winner.  The script is
\texttt{scripts/v2/analyze\_v2.py}; results are in
\texttt{analysis\_v2\_tabular\_ci.csv}
and
\texttt{analysis\_v2\_dl\_ci.csv}.  The complete bootstrap CI table
is shown in Table~\ref{tab:v2_tabular_ci}.

\subsection{Results}
\label{sec:v2-results}

\paragraph{Headline.}  Table~\ref{tab:v2_headline} summarises the
per-family winners selected by validation PR-AUC.  All eight metrics
(accuracy, default-cutoff F1, validation-tuned F1, validation-tuned
precision and recall, ROC-AUC with bootstrap CI, PR-AUC with
bootstrap CI, Brier score) are reported.  The test cohort contains
four patients with two-to-five positive samples per cell, so the
bootstrap intervals on PR-AUC are wide; ROC-AUC, which ranks
predictions rather than counting them, is the more stable quantity.

\begin{table}[t]
\centering
\small
\caption{v2 per-family headline winners on the held-out test set. Tabular rows show the single deterministic fit; the DL row reports mean$\pm$std across five seeds.  All eight metrics shown: accuracy, default-cutoff F1, validation-tuned F1, precision, recall, ROC-AUC (with 95\% bootstrap CI), PR-AUC (with 95\% CI), Brier score.}\label{tab:v2_headline}
\resizebox{\textwidth}{!}{%
\begin{tabular}{llccccccccc}
\toprule
Family & Model & $(T,L,W)$ & Acc & F1 & F1$_{\tau^*}$ & Prec$_{\tau^*}$ & Rec$_{\tau^*}$ & ROC-AUC [95\% CI] & PR-AUC [95\% CI] & Brier \\
\midrule
Tabular & RF & (4,14,14) & 0.988 & 0.000 & 0.667 & 1.000 & 0.500 & 0.934\,[0.802,1.000] & 0.555\,[0.033,1.000] & 0.048 \\
Tabular & XGB & (3,14,14) & 0.972 & 0.308 & 0.067 & 0.035 & 0.750 & 0.875\,[0.633,1.000] & 0.526\,[0.018,1.000] & 0.031 \\
Tabular & LR & (4,14,0) & 0.866 & 0.041 & 0.017 & 0.009 & 0.250 & 0.615\,[0.384,0.919] & 0.022\,[0.006,0.084] & 0.108 \\
Deep & RNN (5 seeds) & (3,14,14) & 0.988 & 0.000 & --- & --- & --- & 0.648 $\pm$ 0.263 & 0.070 $\pm$ 0.076 & 0.012 \\
\bottomrule
\end{tabular}}
\end{table}


\paragraph{Tabular winners across the full grid.}
Table~\ref{tab:v2_tabular_full} reports every metric we record on the
test set for all $81$~$(T,L,W,\text{algo})$ winners.  Accuracy stays
above $0.95$ for almost every cell because the negative class
dominates; the default-cutoff precision/recall/F1 columns are zero
for the majority of cells, which is why we additionally report
metrics at the validation-tuned threshold $\tau^*$.

\begin{longtable}{lcccccccccccc}
\caption{Full tabular results on the held-out test partition for every $(T,L,W,\text{algo})$ winner (81 rows). Metrics at the default $0.5$ cutoff: accuracy, precision, recall, F1.  Metrics at the validation-tuned cutoff $\tau^*$: F1, precision, recall.  Ranking metrics: ROC-AUC, PR-AUC, Brier score.  $n_+$ is the test-positive count.}\label{tab:v2_tabular_full}\\
\toprule
Algo & $(T,L,W)$ & $n$ & $n_+$ & Acc & Prec & Rec & F1 & $\tau^*$ & F1$_{\tau^*}$ & ROC & PR & Brier \\
\midrule
\endfirsthead
\toprule
Algo & $(T,L,W)$ & $n$ & $n_+$ & Acc & Prec & Rec & F1 & $\tau^*$ & F1$_{\tau^*}$ & ROC & PR & Brier \\
\midrule
\endhead
\bottomrule
\endfoot
LR & (2,3,0) & 408 & 3 & 0.831 & 0.000 & 0.000 & 0.000 & 0.000 & 0.017 & 0.464 & 0.010 & 0.143 \\
LR & (2,3,7) & 394 & 3 & 0.794 & 0.000 & 0.000 & 0.000 & 0.000 & 0.008 & 0.347 & 0.007 & 0.168 \\
LR & (2,3,14) & 382 & 3 & 0.668 & 0.000 & 0.000 & 0.000 & 0.000 & 0.017 & 0.382 & 0.008 & 0.304 \\
LR & (2,7,0) & 384 & 3 & 0.745 & 0.000 & 0.000 & 0.000 & 0.000 & 0.000 & 0.206 & 0.006 & 0.218 \\
LR & (2,7,7) & 370 & 3 & 0.868 & 0.000 & 0.000 & 0.000 & 0.000 & 0.000 & 0.338 & 0.008 & 0.114 \\
LR & (2,7,14) & 358 & 3 & 0.553 & 0.000 & 0.000 & 0.000 & 0.000 & 0.009 & 0.348 & 0.008 & 0.435 \\
LR & (2,14,0) & 348 & 2 & 0.899 & 0.000 & 0.000 & 0.000 & 0.001 & 0.000 & 0.483 & 0.009 & 0.087 \\
LR & (2,14,7) & 336 & 2 & 0.714 & 0.000 & 0.000 & 0.000 & 0.051 & 0.000 & 0.313 & 0.006 & 0.251 \\
LR & (2,14,14) & 322 & 2 & 0.748 & 0.000 & 0.000 & 0.000 & 0.278 & 0.000 & 0.355 & 0.007 & 0.212 \\
LR & (3,3,0) & 410 & 5 & 0.751 & 0.000 & 0.000 & 0.000 & 0.000 & 0.015 & 0.265 & 0.010 & 0.204 \\
LR & (3,3,7) & 396 & 5 & 0.816 & 0.000 & 0.000 & 0.000 & 0.000 & 0.007 & 0.201 & 0.009 & 0.165 \\
LR & (3,3,14) & 384 & 5 & 0.630 & 0.000 & 0.000 & 0.000 & 0.000 & 0.016 & 0.237 & 0.010 & 0.355 \\
LR & (3,7,0) & 386 & 5 & 0.725 & 0.000 & 0.000 & 0.000 & 0.000 & 0.008 & 0.261 & 0.010 & 0.254 \\
LR & (3,7,7) & 372 & 5 & 0.742 & 0.000 & 0.000 & 0.000 & 0.000 & 0.032 & 0.464 & 0.014 & 0.243 \\
LR & (3,7,14) & 360 & 5 & 0.578 & 0.000 & 0.000 & 0.000 & 0.000 & 0.017 & 0.366 & 0.012 & 0.406 \\
LR & (3,14,0) & 350 & 4 & 0.889 & 0.000 & 0.000 & 0.000 & 0.002 & 0.000 & 0.533 & 0.015 & 0.086 \\
LR & (3,14,7) & 338 & 4 & 0.840 & 0.000 & 0.000 & 0.000 & 0.071 & 0.000 & 0.385 & 0.012 & 0.131 \\
LR & (3,14,14) & 324 & 4 & 0.833 & 0.000 & 0.000 & 0.000 & 0.627 & 0.000 & 0.376 & 0.012 & 0.146 \\
LR & (4,3,0) & 410 & 5 & 0.912 & 0.000 & 0.000 & 0.000 & 0.000 & 0.015 & 0.381 & 0.012 & 0.079 \\
LR & (4,3,7) & 397 & 5 & 0.909 & 0.000 & 0.000 & 0.000 & 0.000 & 0.008 & 0.255 & 0.010 & 0.085 \\
LR & (4,3,14) & 385 & 5 & 0.732 & 0.000 & 0.000 & 0.000 & 0.000 & 0.015 & 0.295 & 0.011 & 0.224 \\
LR & (4,7,0) & 386 & 5 & 0.741 & 0.000 & 0.000 & 0.000 & 0.000 & 0.008 & 0.285 & 0.010 & 0.226 \\
LR & (4,7,7) & 373 & 5 & 0.783 & 0.000 & 0.000 & 0.000 & 0.000 & 0.017 & 0.436 & 0.013 & 0.203 \\
LR & (4,7,14) & 361 & 5 & 0.618 & 0.000 & 0.000 & 0.000 & 0.047 & 0.025 & 0.424 & 0.014 & 0.246 \\
LR & (4,14,0) & 350 & 4 & 0.866 & 0.022 & 0.250 & 0.041 & 0.001 & 0.017 & 0.615 & 0.022 & 0.108 \\
LR & (4,14,7) & 339 & 4 & 0.832 & 0.018 & 0.250 & 0.034 & 0.204 & 0.026 & 0.541 & 0.018 & 0.140 \\
LR & (4,14,14) & 325 & 4 & 0.938 & 0.000 & 0.000 & 0.000 & 0.246 & 0.000 & 0.509 & 0.021 & 0.059 \\
RF & (2,3,0) & 408 & 3 & 0.993 & 0.000 & 0.000 & 0.000 & 0.179 & 0.000 & 0.531 & 0.012 & 0.023 \\
RF & (2,3,7) & 394 & 3 & 0.992 & 0.000 & 0.000 & 0.000 & 0.273 & 0.029 & 0.574 & 0.017 & 0.051 \\
RF & (2,3,14) & 382 & 3 & 0.992 & 0.000 & 0.000 & 0.000 & 0.212 & 0.021 & 0.696 & 0.030 & 0.054 \\
RF & (2,7,0) & 384 & 3 & 0.969 & 0.000 & 0.000 & 0.000 & 0.235 & 0.026 & 0.499 & 0.039 & 0.043 \\
RF & (2,7,7) & 370 & 3 & 0.992 & 0.000 & 0.000 & 0.000 & 0.180 & 0.017 & 0.439 & 0.039 & 0.034 \\
RF & (2,7,14) & 358 & 3 & 0.966 & 0.000 & 0.000 & 0.000 & 0.278 & 0.014 & 0.536 & 0.040 & 0.078 \\
RF & (2,14,0) & 348 & 2 & 0.994 & 0.000 & 0.000 & 0.000 & 0.487 & 0.000 & 0.828 & 0.024 & 0.038 \\
RF & (2,14,7) & 336 & 2 & 0.994 & 0.000 & 0.000 & 0.000 & 0.391 & 0.000 & 0.614 & 0.011 & 0.033 \\
RF & (2,14,14) & 322 & 2 & 0.994 & 0.000 & 0.000 & 0.000 & 0.413 & 0.000 & 0.841 & 0.510 & 0.051 \\
RF & (3,3,0) & 410 & 5 & 0.988 & 0.000 & 0.000 & 0.000 & 0.171 & 0.028 & 0.499 & 0.017 & 0.030 \\
RF & (3,3,7) & 396 & 5 & 0.987 & 0.000 & 0.000 & 0.000 & 0.151 & 0.022 & 0.430 & 0.014 & 0.036 \\
RF & (3,3,14) & 384 & 5 & 0.987 & 0.000 & 0.000 & 0.000 & 0.209 & 0.027 & 0.532 & 0.026 & 0.061 \\
RF & (3,7,0) & 386 & 5 & 0.964 & 0.000 & 0.000 & 0.000 & 0.266 & 0.062 & 0.557 & 0.033 & 0.044 \\
RF & (3,7,7) & 372 & 5 & 0.981 & 0.000 & 0.000 & 0.000 & 0.225 & 0.069 & 0.513 & 0.023 & 0.034 \\
RF & (3,7,14) & 360 & 5 & 0.981 & 0.000 & 0.000 & 0.000 & 0.218 & 0.030 & 0.534 & 0.033 & 0.041 \\
RF & (3,14,0) & 350 & 4 & 0.989 & 0.000 & 0.000 & 0.000 & 0.443 & 0.000 & 0.788 & 0.033 & 0.028 \\
RF & (3,14,7) & 338 & 4 & 0.988 & 0.000 & 0.000 & 0.000 & 0.416 & 0.000 & 0.818 & 0.042 & 0.049 \\
RF & (3,14,14) & 324 & 4 & 0.988 & 0.000 & 0.000 & 0.000 & 0.418 & 0.000 & 0.916 & 0.554 & 0.050 \\
RF & (4,3,0) & 410 & 5 & 0.988 & 0.000 & 0.000 & 0.000 & 0.210 & 0.000 & 0.637 & 0.025 & 0.028 \\
RF & (4,3,7) & 397 & 5 & 0.987 & 0.000 & 0.000 & 0.000 & 0.179 & 0.019 & 0.521 & 0.017 & 0.037 \\
RF & (4,3,14) & 385 & 5 & 0.987 & 0.000 & 0.000 & 0.000 & 0.177 & 0.030 & 0.661 & 0.033 & 0.045 \\
RF & (4,7,0) & 386 & 5 & 0.987 & 0.000 & 0.000 & 0.000 & 0.262 & 0.094 & 0.670 & 0.054 & 0.049 \\
RF & (4,7,7) & 373 & 5 & 0.987 & 0.000 & 0.000 & 0.000 & 0.149 & 0.034 & 0.576 & 0.043 & 0.029 \\
RF & (4,7,14) & 361 & 5 & 0.981 & 0.000 & 0.000 & 0.000 & 0.190 & 0.051 & 0.708 & 0.064 & 0.042 \\
RF & (4,14,0) & 350 & 4 & 0.989 & 0.000 & 0.000 & 0.000 & 0.428 & 0.000 & 0.818 & 0.316 & 0.055 \\
RF & (4,14,7) & 339 & 4 & 0.988 & 0.000 & 0.000 & 0.000 & 0.360 & 0.000 & 0.879 & 0.327 & 0.027 \\
RF & (4,14,14) & 325 & 4 & 0.988 & 0.000 & 0.000 & 0.000 & 0.354 & 0.667 & 0.934 & 0.555 & 0.048 \\
XGB & (2,3,0) & 408 & 3 & 0.990 & 0.000 & 0.000 & 0.000 & 0.011 & 0.018 & 0.612 & 0.015 & 0.012 \\
XGB & (2,3,7) & 394 & 3 & 0.990 & 0.000 & 0.000 & 0.000 & 0.006 & 0.016 & 0.572 & 0.020 & 0.019 \\
XGB & (2,3,14) & 382 & 3 & 0.969 & 0.000 & 0.000 & 0.000 & 0.003 & 0.016 & 0.653 & 0.032 & 0.031 \\
XGB & (2,7,0) & 384 & 3 & 0.878 & 0.000 & 0.000 & 0.000 & 0.017 & 0.013 & 0.459 & 0.011 & 0.094 \\
XGB & (2,7,7) & 370 & 3 & 0.919 & 0.034 & 0.333 & 0.062 & 0.008 & 0.018 & 0.668 & 0.051 & 0.075 \\
XGB & (2,7,14) & 358 & 3 & 0.913 & 0.033 & 0.333 & 0.061 & 0.001 & 0.017 & 0.638 & 0.352 & 0.054 \\
XGB & (2,14,0) & 348 & 2 & 0.983 & 0.000 & 0.000 & 0.000 & 0.174 & 0.034 & 0.845 & 0.072 & 0.023 \\
XGB & (2,14,7) & 336 & 2 & 0.964 & 0.083 & 0.500 & 0.143 & 0.146 & 0.022 & 0.844 & 0.176 & 0.035 \\
XGB & (2,14,14) & 322 & 2 & 0.904 & 0.032 & 0.500 & 0.061 & 0.115 & 0.020 & 0.659 & 0.505 & 0.074 \\
XGB & (3,3,0) & 410 & 5 & 0.988 & 0.000 & 0.000 & 0.000 & 0.009 & 0.017 & 0.388 & 0.013 & 0.021 \\
XGB & (3,3,7) & 396 & 5 & 0.985 & 0.000 & 0.000 & 0.000 & 0.004 & 0.022 & 0.452 & 0.017 & 0.023 \\
XGB & (3,3,14) & 384 & 5 & 0.979 & 0.000 & 0.000 & 0.000 & 0.002 & 0.026 & 0.541 & 0.019 & 0.031 \\
XGB & (3,7,0) & 386 & 5 & 0.925 & 0.000 & 0.000 & 0.000 & 0.019 & 0.030 & 0.435 & 0.014 & 0.061 \\
XGB & (3,7,7) & 372 & 5 & 0.903 & 0.000 & 0.000 & 0.000 & 0.003 & 0.032 & 0.460 & 0.014 & 0.072 \\
XGB & (3,7,14) & 360 & 5 & 0.956 & 0.077 & 0.200 & 0.111 & 0.004 & 0.030 & 0.713 & 0.050 & 0.034 \\
XGB & (3,14,0) & 350 & 4 & 0.974 & 0.000 & 0.000 & 0.000 & 0.759 & 0.000 & 0.762 & 0.029 & 0.027 \\
XGB & (3,14,7) & 338 & 4 & 0.988 & 0.500 & 0.250 & 0.333 & 0.814 & 0.000 & 0.936 & 0.387 & 0.014 \\
XGB & (3,14,14) & 324 & 4 & 0.972 & 0.222 & 0.500 & 0.308 & 0.120 & 0.067 & 0.875 & 0.526 & 0.031 \\
XGB & (4,3,0) & 410 & 5 & 0.905 & 0.000 & 0.000 & 0.000 & 0.022 & 0.022 & 0.389 & 0.011 & 0.079 \\
XGB & (4,3,7) & 397 & 5 & 0.985 & 0.000 & 0.000 & 0.000 & 0.006 & 0.022 & 0.458 & 0.014 & 0.021 \\
XGB & (4,3,14) & 385 & 5 & 0.966 & 0.000 & 0.000 & 0.000 & 0.001 & 0.026 & 0.689 & 0.050 & 0.037 \\
XGB & (4,7,0) & 386 & 5 & 0.956 & 0.000 & 0.000 & 0.000 & 0.046 & 0.025 & 0.507 & 0.019 & 0.047 \\
XGB & (4,7,7) & 373 & 5 & 0.957 & 0.077 & 0.200 & 0.111 & 0.040 & 0.021 & 0.451 & 0.026 & 0.041 \\
XGB & (4,7,14) & 361 & 5 & 0.967 & 0.182 & 0.400 & 0.250 & 0.005 & 0.022 & 0.770 & 0.133 & 0.041 \\
XGB & (4,14,0) & 350 & 4 & 0.986 & 0.333 & 0.250 & 0.286 & 0.649 & 0.400 & 0.759 & 0.269 & 0.019 \\
XGB & (4,14,7) & 339 & 4 & 0.988 & 0.500 & 0.250 & 0.333 & 0.446 & 0.286 & 0.934 & 0.329 & 0.013 \\
XGB & (4,14,14) & 325 & 4 & 0.938 & 0.100 & 0.500 & 0.167 & 0.394 & 0.133 & 0.843 & 0.296 & 0.048 \\
\end{longtable}


\paragraph{Bootstrap intervals.}
Table~\ref{tab:v2_tabular_ci} attaches a $2{,}000$-resample $95\%$
bootstrap CI to every test ROC-AUC and PR-AUC point estimate.  The
PR-AUC interval reaches the upper bound $1.0$ in every winning
configuration, which is the direct consequence of having
$n_{\text{test}}^+ \in [2,5]$.

\begin{longtable}{lccccc}
\caption{Bootstrap $95\%$ confidence intervals (2,000 resamples) on test PR-AUC and ROC-AUC for every tabular winner.}\label{tab:v2_tabular_ci}\\
\toprule
Algo & $(T,L,W)$ & $n_+$ & PR-AUC [95\% CI] & ROC-AUC [95\% CI] \\
\midrule
\endfirsthead
\toprule
Algo & $(T,L,W)$ & $n_+$ & PR-AUC [95\% CI] & ROC-AUC [95\% CI] \\
\midrule
\endhead
\bottomrule
\endfoot
LR & (2,3,0) & 3 & 0.010\,[0.003,0.028] & 0.464\,[0.076,0.725] \\
LR & (2,3,7) & 3 & 0.007\,[0.003,0.019] & 0.347\,[0.204,0.548] \\
LR & (2,3,14) & 3 & 0.008\,[0.003,0.021] & 0.382\,[0.154,0.536] \\
LR & (2,7,0) & 3 & 0.006\,[0.003,0.017] & 0.206\,[0.016,0.334] \\
LR & (2,7,7) & 3 & 0.008\,[0.004,0.020] & 0.338\,[0.276,0.412] \\
LR & (2,7,14) & 3 & 0.008\,[0.004,0.021] & 0.348\,[0.199,0.496] \\
LR & (2,14,0) & 2 & 0.009\,[0.004,0.030] & 0.483\,[0.226,0.743] \\
LR & (2,14,7) & 2 & 0.006\,[0.003,0.019] & 0.313\,[0.107,0.534] \\
LR & (2,14,14) & 2 & 0.007\,[0.004,0.021] & 0.355\,[0.262,0.449] \\
LR & (3,3,0) & 5 & 0.010\,[0.003,0.023] & 0.265\,[0.048,0.550] \\
LR & (3,3,7) & 5 & 0.009\,[0.004,0.019] & 0.201\,[0.069,0.336] \\
LR & (3,3,14) & 5 & 0.010\,[0.004,0.022] & 0.237\,[0.042,0.444] \\
LR & (3,7,0) & 5 & 0.010\,[0.004,0.021] & 0.261\,[0.108,0.401] \\
LR & (3,7,7) & 5 & 0.014\,[0.005,0.030] & 0.464\,[0.356,0.570] \\
LR & (3,7,14) & 5 & 0.012\,[0.005,0.026] & 0.366\,[0.251,0.494] \\
LR & (3,14,0) & 4 & 0.015\,[0.006,0.034] & 0.533\,[0.351,0.650] \\
LR & (3,14,7) & 4 & 0.012\,[0.005,0.027] & 0.385\,[0.166,0.509] \\
LR & (3,14,14) & 4 & 0.012\,[0.005,0.027] & 0.376\,[0.226,0.480] \\
LR & (4,3,0) & 5 & 0.012\,[0.003,0.032] & 0.381\,[0.171,0.665] \\
LR & (4,3,7) & 5 & 0.010\,[0.004,0.020] & 0.255\,[0.162,0.365] \\
LR & (4,3,14) & 5 & 0.011\,[0.004,0.025] & 0.295\,[0.131,0.539] \\
LR & (4,7,0) & 5 & 0.010\,[0.004,0.022] & 0.285\,[0.174,0.391] \\
LR & (4,7,7) & 5 & 0.013\,[0.006,0.031] & 0.436\,[0.306,0.597] \\
LR & (4,7,14) & 5 & 0.014\,[0.006,0.032] & 0.424\,[0.193,0.614] \\
LR & (4,14,0) & 4 & 0.022\,[0.006,0.084] & 0.615\,[0.384,0.919] \\
LR & (4,14,7) & 4 & 0.018\,[0.006,0.062] & 0.541\,[0.207,0.855] \\
LR & (4,14,14) & 4 & 0.021\,[0.005,0.083] & 0.509\,[0.306,0.897] \\
RF & (2,3,0) & 3 & 0.012\,[0.003,0.049] & 0.531\,[0.161,0.872] \\
RF & (2,3,7) & 3 & 0.017\,[0.003,0.076] & 0.574\,[0.252,0.923] \\
RF & (2,3,14) & 3 & 0.030\,[0.004,0.149] & 0.696\,[0.339,0.961] \\
RF & (2,7,0) & 3 & 0.039\,[0.003,0.250] & 0.499\,[0.225,0.980] \\
RF & (2,7,7) & 3 & 0.039\,[0.003,0.231] & 0.439\,[0.128,0.981] \\
RF & (2,7,14) & 3 & 0.040\,[0.004,0.240] & 0.536\,[0.247,0.978] \\
RF & (2,14,0) & 2 & 0.024\,[0.012,0.075] & 0.828\,[0.758,0.893] \\
RF & (2,14,7) & 2 & 0.011\,[0.006,0.033] & 0.614\,[0.484,0.737] \\
RF & (2,14,14) & 2 & 0.510\,[0.009,1.000] & 0.841\,[0.647,1.000] \\
RF & (3,3,0) & 5 & 0.017\,[0.005,0.052] & 0.499\,[0.143,0.817] \\
RF & (3,3,7) & 5 & 0.014\,[0.005,0.038] & 0.430\,[0.134,0.712] \\
RF & (3,3,14) & 5 & 0.026\,[0.005,0.118] & 0.532\,[0.173,0.888] \\
RF & (3,7,0) & 5 & 0.033\,[0.006,0.177] & 0.557\,[0.313,0.816] \\
RF & (3,7,7) & 5 & 0.023\,[0.005,0.106] & 0.513\,[0.302,0.762] \\
RF & (3,7,14) & 5 & 0.033\,[0.006,0.177] & 0.534\,[0.303,0.805] \\
RF & (3,14,0) & 4 & 0.033\,[0.013,0.095] & 0.788\,[0.672,0.911] \\
RF & (3,14,7) & 4 & 0.042\,[0.016,0.112] & 0.818\,[0.641,0.914] \\
RF & (3,14,14) & 4 & 0.554\,[0.025,1.000] & 0.916\,[0.727,1.000] \\
RF & (4,3,0) & 5 & 0.025\,[0.006,0.074] & 0.637\,[0.341,0.879] \\
RF & (4,3,7) & 5 & 0.017\,[0.005,0.051] & 0.521\,[0.241,0.787] \\
RF & (4,3,14) & 5 & 0.033\,[0.008,0.117] & 0.661\,[0.352,0.922] \\
RF & (4,7,0) & 5 & 0.054\,[0.007,0.195] & 0.670\,[0.299,0.964] \\
RF & (4,7,7) & 5 & 0.043\,[0.005,0.190] & 0.576\,[0.275,0.953] \\
RF & (4,7,14) & 5 & 0.064\,[0.009,0.254] & 0.708\,[0.398,0.969] \\
RF & (4,14,0) & 4 & 0.316\,[0.011,1.000] & 0.818\,[0.527,1.000] \\
RF & (4,14,7) & 4 & 0.327\,[0.016,1.000] & 0.879\,[0.648,1.000] \\
RF & (4,14,14) & 4 & 0.555\,[0.033,1.000] & 0.934\,[0.802,1.000] \\
XGB & (2,3,0) & 3 & 0.015\,[0.004,0.056] & 0.612\,[0.368,0.894] \\
XGB & (2,3,7) & 3 & 0.020\,[0.004,0.102] & 0.572\,[0.321,0.944] \\
XGB & (2,3,14) & 3 & 0.032\,[0.003,0.161] & 0.653\,[0.217,0.968] \\
XGB & (2,7,0) & 3 & 0.011\,[0.003,0.040] & 0.459\,[0.026,0.825] \\
XGB & (2,7,7) & 3 & 0.051\,[0.003,0.228] & 0.668\,[0.119,0.978] \\
XGB & (2,7,14) & 3 & 0.352\,[0.003,1.000] & 0.638\,[0.022,1.000] \\
XGB & (2,14,0) & 2 & 0.072\,[0.009,0.375] & 0.845\,[0.680,0.988] \\
XGB & (2,14,7) & 2 & 0.176\,[0.009,1.000] & 0.844\,[0.661,1.000] \\
XGB & (2,14,14) & 2 & 0.505\,[0.004,1.000] & 0.659\,[0.287,1.000] \\
XGB & (3,3,0) & 5 & 0.013\,[0.004,0.044] & 0.388\,[0.081,0.744] \\
XGB & (3,3,7) & 5 & 0.017\,[0.005,0.058] & 0.452\,[0.087,0.847] \\
XGB & (3,3,14) & 5 & 0.019\,[0.005,0.052] & 0.541\,[0.173,0.805] \\
XGB & (3,7,0) & 5 & 0.014\,[0.004,0.039] & 0.435\,[0.247,0.688] \\
XGB & (3,7,7) & 5 & 0.014\,[0.005,0.035] & 0.460\,[0.259,0.661] \\
XGB & (3,7,14) & 5 & 0.050\,[0.012,0.224] & 0.713\,[0.469,0.922] \\
XGB & (3,14,0) & 4 & 0.029\,[0.012,0.071] & 0.762\,[0.628,0.861] \\
XGB & (3,14,7) & 4 & 0.387\,[0.025,1.000] & 0.936\,[0.782,1.000] \\
XGB & (3,14,14) & 4 & 0.526\,[0.018,1.000] & 0.875\,[0.633,1.000] \\
XGB & (4,3,0) & 5 & 0.011\,[0.004,0.026] & 0.389\,[0.196,0.586] \\
XGB & (4,3,7) & 5 & 0.014\,[0.006,0.035] & 0.458\,[0.154,0.690] \\
XGB & (4,3,14) & 5 & 0.050\,[0.009,0.201] & 0.689\,[0.385,0.963] \\
XGB & (4,7,0) & 5 & 0.019\,[0.005,0.061] & 0.507\,[0.184,0.849] \\
XGB & (4,7,7) & 5 & 0.026\,[0.005,0.142] & 0.451\,[0.176,0.762] \\
XGB & (4,7,14) & 5 & 0.133\,[0.029,0.367] & 0.770\,[0.321,0.982] \\
XGB & (4,14,0) & 4 & 0.269\,[0.009,1.000] & 0.759\,[0.625,1.000] \\
XGB & (4,14,7) & 4 & 0.329\,[0.028,0.986] & 0.934\,[0.870,0.999] \\
XGB & (4,14,14) & 4 & 0.296\,[0.014,1.000] & 0.843\,[0.651,1.000] \\
\end{longtable}


\paragraph{Deep-learning seed sensitivity.}
Table~\ref{tab:v2_dl_seeds} lists every metric for each of the five
seeds of the RNN winner at $(T,L,W)\!=\!(3,14,14)$.  The default-cutoff
F1 is zero across all five seeds (the network produces well-ranked
but small-magnitude probabilities), but the ROC-AUC mean of $0.902$
with std $0.057$ confirms that the ranking signal is reproducible
across initialisations.

\begin{table}[t]
\centering
\small
\caption{Deep-learning winner (RNN, $(T,L,W)=(3,14,14)$, all sensors) refit across five seeds.  Default-cutoff accuracy / precision / recall / F1 plus ranking metrics on the held-out test partition.}\label{tab:v2_dl_seeds}
\begin{tabular}{ccccccccc}
\toprule
Seed & Best Epoch & Acc & Prec & Rec & F1 & ROC-AUC & PR-AUC & Brier \\
\midrule
42 & 27 & 0.988 & 0.000 & 0.000 & 0.000 & 0.962 & 0.357 & 0.012 \\
43 & 28 & 0.988 & 0.000 & 0.000 & 0.000 & 0.958 & 0.405 & 0.012 \\
44 & 30 & 0.988 & 0.000 & 0.000 & 0.000 & 0.828 & 0.075 & 0.012 \\
45 & 30 & 0.988 & 0.000 & 0.000 & 0.000 & 0.887 & 0.341 & 0.012 \\
46 & 29 & 0.988 & 0.000 & 0.000 & 0.000 & 0.876 & 0.129 & 0.012 \\
\midrule
\textbf{Mean} & --- & 0.988 & 0.000 & 0.000 & 0.000 & 0.902 & 0.261 & 0.012 \\
\textbf{Std} & --- & 0.000 & 0.000 & 0.000 & 0.000 & 0.057 & 0.148 & 0.000 \\
\bottomrule
\end{tabular}
\end{table}


\paragraph{Sensor ablation -- single most striking finding.}
The clearest signal in the v2 experiment is sensor-isolated.  An
XGBoost model trained on \emph{rescue-inhaler features alone} on the
severe-event cell $(T,L,W)\!=\!(4,14,14)$ reaches test
ROC-AUC $=\!0.964$ with PR-AUC $=\!0.164$ and tuned-threshold
F1 $=\!0.136$ (top row of Table~\ref{tab:v2_ablation_topk}, full
record in Table~\ref{tab:v2_sensor_full}).  Combined with the strong
all-sensor RF result on the same cell, this suggests that the
predictive signal for severe events is dominated by rescue-inhaler
usage rather than by smartwatch physiology in this cohort.
Leave-one-out experiments show that removing the
\texttt{smartinhaler} channel drops every algorithm's validation
PR-AUC, while removing other sensors typically does not.

\begin{table}[t]
\centering
\small
\caption{Top-3 sensor subsets per tabular algorithm, ranked by validation PR-AUC.  All metrics shown.}\label{tab:v2_ablation_topk}
\begin{tabular}{llccccccc}
\toprule
Algo & Subset & $(T,L,W)$ & Val PR & Val ROC & Val F1 & Test PR & Test ROC & Test F1 \\
\midrule
LR & smartinhaler\_only & (3,14,0) & 0.036 & 0.891 & 0.069 & 0.018 & 0.596 & 0.040 \\
LR & smartinhaler\_only & (4,14,0) & 0.036 & 0.891 & 0.069 & 0.024 & 0.671 & 0.033 \\
LR & smartinhaler\_only & (2,14,0) & 0.036 & 0.891 & 0.069 & 0.008 & 0.465 & 0.000 \\
RF & no\_smartinhaler & (3,14,7) & 1.000 & 1.000 & 0.400 & 0.012 & 0.396 & 0.000 \\
RF & all & (4,14,7) & 0.250 & 0.989 & 0.000 & 0.031 & 0.664 & 0.000 \\
RF & all & (3,14,0) & 0.250 & 0.984 & 0.000 & 0.011 & 0.374 & 0.000 \\
XGB & smartinhaler\_only & (4,14,14) & 0.067 & 0.962 & 0.111 & 0.164 & 0.964 & 0.136 \\
XGB & smartinhaler\_only & (2,14,7) & 0.056 & 0.945 & 0.000 & 0.014 & 0.696 & 0.000 \\
XGB & smartinhaler\_only & (4,14,7) & 0.053 & 0.940 & 0.000 & 0.059 & 0.746 & 0.061 \\
\bottomrule
\end{tabular}
\end{table}


\begin{longtable}{llcccccccccc}
\caption{Full sensor-ablation grid: every $(T,L,W,\text{algo},\text{subset})$ refit.  Columns report validation and test PR-AUC / ROC-AUC / F1 plus the test accuracy.  729 rows.}\label{tab:v2_sensor_full}\\
\toprule
Algo & Subset & $(T,L,W)$ & $n_{\text{feat}}$ & Val PR & Val ROC & Val F1 & Test Acc & Test Prec & Test Rec & Test ROC & Test PR \\
\midrule
\endfirsthead
\toprule
Algo & Subset & $(T,L,W)$ & $n_{\text{feat}}$ & Val PR & Val ROC & Val F1 & Test Acc & Test Prec & Test Rec & Test ROC & Test PR \\
\midrule
\endhead
\bottomrule
\endfoot
LR & all & (2,3,0) & 427 & 0.006 & 0.255 & 0.000 & 0.821 & 0.014 & 0.333 & 0.471 & 0.012 \\
LR & environment\_only & (2,3,0) & 171 & 0.005 & 0.162 & 0.000 & 0.897 & 0.000 & 0.000 & 0.297 & 0.007 \\
LR & no\_environment & (2,3,0) & 315 & 0.006 & 0.287 & 0.000 & 0.833 & 0.000 & 0.000 & 0.439 & 0.010 \\
LR & no\_peakflow & (2,3,0) & 355 & 0.006 & 0.282 & 0.000 & 0.946 & 0.048 & 0.333 & 0.694 & 0.052 \\
LR & no\_smartinhaler & (2,3,0) & 400 & 0.007 & 0.310 & 0.000 & 0.836 & 0.000 & 0.000 & 0.456 & 0.011 \\
LR & no\_smartwatch & (2,3,0) & 270 & 0.005 & 0.144 & 0.000 & 0.811 & 0.000 & 0.000 & 0.205 & 0.006 \\
LR & peakflow\_only & (2,3,0) & 131 & 0.006 & 0.167 & 0.000 & 0.674 & 0.000 & 0.000 & 0.142 & 0.005 \\
LR & smartinhaler\_only & (2,3,0) & 86 & 0.006 & 0.282 & 0.000 & 0.990 & 0.000 & 0.000 & 0.120 & 0.005 \\
LR & smartwatch\_only & (2,3,0) & 216 & 0.008 & 0.426 & 0.000 & 0.941 & 0.000 & 0.000 & 0.686 & 0.018 \\
LR & all & (2,3,7) & 427 & 0.006 & 0.181 & 0.000 & 0.802 & 0.000 & 0.000 & 0.246 & 0.006 \\
LR & environment\_only & (2,3,7) & 171 & 0.005 & 0.088 & 0.000 & 0.865 & 0.000 & 0.000 & 0.282 & 0.007 \\
LR & no\_environment & (2,3,7) & 315 & 0.006 & 0.287 & 0.000 & 0.690 & 0.000 & 0.000 & 0.215 & 0.006 \\
LR & no\_peakflow & (2,3,7) & 355 & 0.006 & 0.227 & 0.000 & 0.924 & 0.000 & 0.000 & 0.526 & 0.011 \\
LR & no\_smartinhaler & (2,3,7) & 400 & 0.006 & 0.190 & 0.000 & 0.807 & 0.000 & 0.000 & 0.289 & 0.007 \\
LR & no\_smartwatch & (2,3,7) & 270 & 0.005 & 0.130 & 0.000 & 0.792 & 0.000 & 0.000 & 0.171 & 0.006 \\
LR & peakflow\_only & (2,3,7) & 131 & 0.005 & 0.104 & 0.000 & 0.541 & 0.000 & 0.000 & 0.120 & 0.006 \\
LR & smartinhaler\_only & (2,3,7) & 86 & 0.005 & 0.213 & 0.000 & 0.980 & 0.000 & 0.000 & 0.077 & 0.005 \\
LR & smartwatch\_only & (2,3,7) & 216 & 0.008 & 0.417 & 0.000 & 0.891 & 0.000 & 0.000 & 0.613 & 0.013 \\
LR & all & (2,3,14) & 427 & 0.005 & 0.157 & 0.000 & 0.686 & 0.000 & 0.000 & 0.317 & 0.007 \\
LR & environment\_only & (2,3,14) & 171 & 0.005 & 0.176 & 0.000 & 0.809 & 0.000 & 0.000 & 0.354 & 0.008 \\
LR & no\_environment & (2,3,14) & 315 & 0.006 & 0.255 & 0.000 & 0.647 & 0.000 & 0.000 & 0.201 & 0.006 \\
LR & no\_peakflow & (2,3,14) & 355 & 0.005 & 0.171 & 0.000 & 0.901 & 0.000 & 0.000 & 0.362 & 0.008 \\
LR & no\_smartinhaler & (2,3,14) & 400 & 0.006 & 0.185 & 0.000 & 0.683 & 0.000 & 0.000 & 0.303 & 0.007 \\
LR & no\_smartwatch & (2,3,14) & 270 & 0.005 & 0.144 & 0.000 & 0.639 & 0.000 & 0.000 & 0.215 & 0.006 \\
LR & peakflow\_only & (2,3,14) & 131 & 0.005 & 0.109 & 0.000 & 0.406 & 0.000 & 0.000 & 0.153 & 0.006 \\
LR & smartinhaler\_only & (2,3,14) & 86 & 0.005 & 0.005 & 0.000 & 0.864 & 0.000 & 0.000 & 0.066 & 0.005 \\
LR & smartwatch\_only & (2,3,14) & 216 & 0.006 & 0.310 & 0.000 & 0.916 & 0.000 & 0.000 & 0.514 & 0.011 \\
LR & all & (2,7,0) & 427 & 0.006 & 0.208 & 0.000 & 0.820 & 0.000 & 0.000 & 0.167 & 0.006 \\
LR & environment\_only & (2,7,0) & 171 & 0.005 & 0.045 & 0.000 & 0.776 & 0.012 & 0.333 & 0.408 & 0.010 \\
LR & no\_environment & (2,7,0) & 315 & 0.006 & 0.218 & 0.000 & 0.518 & 0.000 & 0.000 & 0.132 & 0.006 \\
LR & no\_peakflow & (2,7,0) & 355 & 0.006 & 0.139 & 0.000 & 0.896 & 0.000 & 0.000 & 0.244 & 0.007 \\
LR & no\_smartinhaler & (2,7,0) & 400 & 0.006 & 0.223 & 0.000 & 0.789 & 0.000 & 0.000 & 0.169 & 0.006 \\
LR & no\_smartwatch & (2,7,0) & 270 & 0.006 & 0.114 & 0.000 & 0.654 & 0.008 & 0.333 & 0.354 & 0.008 \\
LR & peakflow\_only & (2,7,0) & 131 & 0.006 & 0.238 & 0.000 & 0.547 & 0.000 & 0.000 & 0.191 & 0.006 \\
LR & smartinhaler\_only & (2,7,0) & 86 & 0.006 & 0.178 & 0.000 & 0.992 & 0.000 & 0.000 & 0.325 & 0.009 \\
LR & smartwatch\_only & (2,7,0) & 216 & 0.008 & 0.376 & 0.000 & 0.737 & 0.000 & 0.000 & 0.180 & 0.006 \\
LR & all & (2,7,7) & 427 & 0.006 & 0.163 & 0.000 & 0.927 & 0.000 & 0.000 & 0.354 & 0.008 \\
LR & environment\_only & (2,7,7) & 171 & 0.005 & 0.015 & 0.000 & 0.749 & 0.011 & 0.333 & 0.404 & 0.013 \\
LR & no\_environment & (2,7,7) & 315 & 0.006 & 0.228 & 0.000 & 0.570 & 0.000 & 0.000 & 0.153 & 0.006 \\
LR & no\_peakflow & (2,7,7) & 355 & 0.006 & 0.124 & 0.000 & 0.935 & 0.043 & 0.333 & 0.518 & 0.025 \\
LR & no\_smartinhaler & (2,7,7) & 400 & 0.006 & 0.144 & 0.000 & 0.930 & 0.000 & 0.000 & 0.268 & 0.007 \\
LR & no\_smartwatch & (2,7,7) & 270 & 0.005 & 0.109 & 0.000 & 0.727 & 0.010 & 0.333 & 0.450 & 0.047 \\
LR & peakflow\_only & (2,7,7) & 131 & 0.006 & 0.188 & 0.000 & 0.551 & 0.000 & 0.000 & 0.189 & 0.006 \\
LR & smartinhaler\_only & (2,7,7) & 86 & 0.006 & 0.265 & 0.000 & 0.973 & 0.000 & 0.000 & 0.591 & 0.033 \\
LR & smartwatch\_only & (2,7,7) & 216 & 0.006 & 0.233 & 0.000 & 0.765 & 0.000 & 0.000 & 0.169 & 0.006 \\
LR & all & (2,7,14) & 427 & 0.005 & 0.064 & 0.000 & 0.422 & 0.005 & 0.333 & 0.344 & 0.008 \\
LR & environment\_only & (2,7,14) & 171 & 0.005 & 0.015 & 0.000 & 0.656 & 0.008 & 0.333 & 0.530 & 0.011 \\
LR & no\_environment & (2,7,14) & 315 & 0.006 & 0.252 & 0.000 & 0.458 & 0.000 & 0.000 & 0.153 & 0.007 \\
LR & no\_peakflow & (2,7,14) & 355 & 0.005 & 0.064 & 0.000 & 0.791 & 0.014 & 0.333 & 0.631 & 0.018 \\
LR & no\_smartinhaler & (2,7,14) & 400 & 0.005 & 0.079 & 0.000 & 0.453 & 0.000 & 0.000 & 0.223 & 0.007 \\
LR & no\_smartwatch & (2,7,14) & 270 & 0.005 & 0.015 & 0.000 & 0.316 & 0.004 & 0.333 & 0.380 & 0.011 \\
LR & peakflow\_only & (2,7,14) & 131 & 0.006 & 0.238 & 0.000 & 0.397 & 0.000 & 0.000 & 0.186 & 0.007 \\
LR & smartinhaler\_only & (2,7,14) & 86 & 0.005 & 0.072 & 0.000 & 0.807 & 0.015 & 0.333 & 0.372 & 0.116 \\
LR & smartwatch\_only & (2,7,14) & 216 & 0.008 & 0.396 & 0.000 & 0.760 & 0.000 & 0.000 & 0.171 & 0.007 \\
LR & all & (2,14,0) & 427 & 0.007 & 0.251 & 0.000 & 0.828 & 0.000 & 0.000 & 0.616 & 0.011 \\
LR & environment\_only & (2,14,0) & 171 & 0.007 & 0.257 & 0.000 & 0.871 & 0.000 & 0.000 & 0.730 & 0.016 \\
LR & no\_environment & (2,14,0) & 315 & 0.007 & 0.175 & 0.000 & 0.514 & 0.000 & 0.000 & 0.188 & 0.005 \\
LR & no\_peakflow & (2,14,0) & 355 & 0.007 & 0.230 & 0.000 & 0.853 & 0.000 & 0.000 & 0.578 & 0.011 \\
LR & no\_smartinhaler & (2,14,0) & 400 & 0.007 & 0.219 & 0.000 & 0.750 & 0.000 & 0.000 & 0.558 & 0.009 \\
LR & no\_smartwatch & (2,14,0) & 270 & 0.010 & 0.448 & 0.000 & 0.802 & 0.014 & 0.500 & 0.812 & 0.508 \\
LR & peakflow\_only & (2,14,0) & 131 & 0.006 & 0.148 & 0.000 & 0.328 & 0.004 & 0.500 & 0.410 & 0.008 \\
LR & smartinhaler\_only & (2,14,0) & 86 & 0.036 & 0.891 & 0.069 & 0.897 & 0.000 & 0.000 & 0.465 & 0.008 \\
LR & smartwatch\_only & (2,14,0) & 216 & 0.009 & 0.421 & 0.000 & 0.830 & 0.000 & 0.000 & 0.117 & 0.005 \\
LR & all & (2,14,7) & 427 & 0.007 & 0.180 & 0.000 & 0.509 & 0.000 & 0.000 & 0.400 & 0.007 \\
LR & environment\_only & (2,14,7) & 171 & 0.008 & 0.355 & 0.000 & 0.824 & 0.000 & 0.000 & 0.722 & 0.016 \\
LR & no\_environment & (2,14,7) & 315 & 0.008 & 0.306 & 0.000 & 0.482 & 0.000 & 0.000 & 0.225 & 0.006 \\
LR & no\_peakflow & (2,14,7) & 355 & 0.007 & 0.197 & 0.000 & 0.631 & 0.000 & 0.000 & 0.301 & 0.006 \\
LR & no\_smartinhaler & (2,14,7) & 400 & 0.006 & 0.158 & 0.000 & 0.443 & 0.000 & 0.000 & 0.361 & 0.007 \\
LR & no\_smartwatch & (2,14,7) & 270 & 0.009 & 0.421 & 0.000 & 0.673 & 0.009 & 0.500 & 0.704 & 0.023 \\
LR & peakflow\_only & (2,14,7) & 131 & 0.007 & 0.246 & 0.000 & 0.396 & 0.005 & 0.500 & 0.290 & 0.006 \\
LR & smartinhaler\_only & (2,14,7) & 86 & 0.036 & 0.891 & 0.045 & 0.762 & 0.000 & 0.000 & 0.349 & 0.007 \\
LR & smartwatch\_only & (2,14,7) & 216 & 0.008 & 0.279 & 0.000 & 0.640 & 0.000 & 0.000 & 0.081 & 0.005 \\
LR & all & (2,14,14) & 427 & 0.007 & 0.208 & 0.000 & 0.488 & 0.000 & 0.000 & 0.373 & 0.007 \\
LR & environment\_only & (2,14,14) & 171 & 0.009 & 0.426 & 0.000 & 0.786 & 0.000 & 0.000 & 0.628 & 0.012 \\
LR & no\_environment & (2,14,14) & 315 & 0.011 & 0.497 & 0.000 & 0.547 & 0.000 & 0.000 & 0.206 & 0.006 \\
LR & no\_peakflow & (2,14,14) & 355 & 0.007 & 0.251 & 0.000 & 0.736 & 0.000 & 0.000 & 0.406 & 0.008 \\
LR & no\_smartinhaler & (2,14,14) & 400 & 0.007 & 0.180 & 0.000 & 0.475 & 0.000 & 0.000 & 0.231 & 0.006 \\
LR & no\_smartwatch & (2,14,14) & 270 & 0.010 & 0.481 & 0.021 & 0.593 & 0.008 & 0.500 & 0.597 & 0.018 \\
LR & peakflow\_only & (2,14,14) & 131 & 0.009 & 0.443 & 0.000 & 0.398 & 0.005 & 0.500 & 0.269 & 0.006 \\
LR & smartinhaler\_only & (2,14,14) & 86 & 0.008 & 0.355 & 0.000 & 0.739 & 0.012 & 0.500 & 0.572 & 0.011 \\
LR & smartwatch\_only & (2,14,14) & 216 & 0.009 & 0.426 & 0.000 & 0.919 & 0.000 & 0.000 & 0.214 & 0.006 \\
LR & all & (3,3,0) & 427 & 0.006 & 0.222 & 0.000 & 0.766 & 0.000 & 0.000 & 0.260 & 0.010 \\
LR & environment\_only & (3,3,0) & 171 & 0.006 & 0.185 & 0.000 & 0.839 & 0.000 & 0.000 & 0.178 & 0.009 \\
LR & no\_environment & (3,3,0) & 315 & 0.005 & 0.181 & 0.000 & 0.624 & 0.007 & 0.200 & 0.248 & 0.010 \\
LR & no\_peakflow & (3,3,0) & 355 & 0.005 & 0.176 & 0.000 & 0.924 & 0.000 & 0.000 & 0.370 & 0.014 \\
LR & no\_smartinhaler & (3,3,0) & 400 & 0.006 & 0.227 & 0.000 & 0.724 & 0.009 & 0.200 & 0.251 & 0.010 \\
LR & no\_smartwatch & (3,3,0) & 270 & 0.006 & 0.204 & 0.000 & 0.751 & 0.000 & 0.000 & 0.144 & 0.008 \\
LR & peakflow\_only & (3,3,0) & 131 & 0.005 & 0.118 & 0.000 & 0.634 & 0.000 & 0.000 & 0.136 & 0.008 \\
LR & smartinhaler\_only & (3,3,0) & 86 & 0.005 & 0.079 & 0.000 & 0.959 & 0.000 & 0.000 & 0.119 & 0.008 \\
LR & smartwatch\_only & (3,3,0) & 216 & 0.006 & 0.227 & 0.000 & 0.871 & 0.020 & 0.200 & 0.354 & 0.014 \\
LR & all & (3,3,7) & 427 & 0.006 & 0.181 & 0.000 & 0.785 & 0.000 & 0.000 & 0.151 & 0.009 \\
LR & environment\_only & (3,3,7) & 171 & 0.005 & 0.065 & 0.000 & 0.949 & 0.000 & 0.000 & 0.172 & 0.009 \\
LR & no\_environment & (3,3,7) & 315 & 0.006 & 0.213 & 0.000 & 0.694 & 0.000 & 0.000 & 0.177 & 0.009 \\
LR & no\_peakflow & (3,3,7) & 355 & 0.005 & 0.139 & 0.000 & 0.889 & 0.000 & 0.000 & 0.192 & 0.009 \\
LR & no\_smartinhaler & (3,3,7) & 400 & 0.005 & 0.148 & 0.000 & 0.775 & 0.000 & 0.000 & 0.167 & 0.009 \\
LR & no\_smartwatch & (3,3,7) & 270 & 0.005 & 0.116 & 0.000 & 0.818 & 0.000 & 0.000 & 0.121 & 0.008 \\
LR & peakflow\_only & (3,3,7) & 131 & 0.005 & 0.118 & 0.000 & 0.737 & 0.000 & 0.000 & 0.128 & 0.008 \\
LR & smartinhaler\_only & (3,3,7) & 86 & 0.005 & 0.120 & 0.000 & 0.987 & 0.000 & 0.000 & 0.184 & 0.009 \\
LR & smartwatch\_only & (3,3,7) & 216 & 0.006 & 0.204 & 0.000 & 0.854 & 0.000 & 0.000 & 0.337 & 0.011 \\
LR & all & (3,3,14) & 427 & 0.006 & 0.218 & 0.000 & 0.591 & 0.000 & 0.000 & 0.185 & 0.009 \\
LR & environment\_only & (3,3,14) & 171 & 0.005 & 0.153 & 0.000 & 0.794 & 0.000 & 0.000 & 0.174 & 0.009 \\
LR & no\_environment & (3,3,14) & 315 & 0.006 & 0.190 & 0.000 & 0.529 & 0.000 & 0.000 & 0.175 & 0.009 \\
LR & no\_peakflow & (3,3,14) & 355 & 0.005 & 0.176 & 0.000 & 0.763 & 0.000 & 0.000 & 0.217 & 0.010 \\
LR & no\_smartinhaler & (3,3,14) & 400 & 0.005 & 0.176 & 0.000 & 0.589 & 0.000 & 0.000 & 0.180 & 0.009 \\
LR & no\_smartwatch & (3,3,14) & 270 & 0.005 & 0.148 & 0.000 & 0.617 & 0.000 & 0.000 & 0.134 & 0.009 \\
LR & peakflow\_only & (3,3,14) & 131 & 0.006 & 0.181 & 0.000 & 0.451 & 0.000 & 0.000 & 0.106 & 0.008 \\
LR & smartinhaler\_only & (3,3,14) & 86 & 0.005 & 0.157 & 0.000 & 0.862 & 0.000 & 0.000 & 0.102 & 0.008 \\
LR & smartwatch\_only & (3,3,14) & 216 & 0.006 & 0.204 & 0.000 & 0.802 & 0.014 & 0.200 & 0.281 & 0.011 \\
LR & all & (3,7,0) & 427 & 0.007 & 0.257 & 0.000 & 0.834 & 0.000 & 0.000 & 0.188 & 0.009 \\
LR & environment\_only & (3,7,0) & 171 & 0.005 & 0.040 & 0.000 & 0.808 & 0.000 & 0.000 & 0.307 & 0.011 \\
LR & no\_environment & (3,7,0) & 315 & 0.006 & 0.233 & 0.000 & 0.513 & 0.000 & 0.000 & 0.187 & 0.009 \\
LR & no\_peakflow & (3,7,0) & 355 & 0.006 & 0.129 & 0.000 & 0.829 & 0.000 & 0.000 & 0.152 & 0.009 \\
LR & no\_smartinhaler & (3,7,0) & 400 & 0.006 & 0.208 & 0.000 & 0.845 & 0.000 & 0.000 & 0.175 & 0.009 \\
LR & no\_smartwatch & (3,7,0) & 270 & 0.006 & 0.223 & 0.000 & 0.868 & 0.000 & 0.000 & 0.361 & 0.012 \\
LR & peakflow\_only & (3,7,0) & 131 & 0.006 & 0.233 & 0.000 & 0.655 & 0.000 & 0.000 & 0.166 & 0.009 \\
LR & smartinhaler\_only & (3,7,0) & 86 & 0.005 & 0.116 & 0.000 & 0.959 & 0.000 & 0.000 & 0.239 & 0.010 \\
LR & smartwatch\_only & (3,7,0) & 216 & 0.012 & 0.584 & 0.000 & 0.738 & 0.000 & 0.000 & 0.165 & 0.009 \\
LR & all & (3,7,7) & 427 & 0.006 & 0.198 & 0.000 & 0.817 & 0.000 & 0.000 & 0.366 & 0.012 \\
LR & environment\_only & (3,7,7) & 171 & 0.005 & 0.010 & 0.000 & 0.677 & 0.009 & 0.200 & 0.281 & 0.012 \\
LR & no\_environment & (3,7,7) & 315 & 0.006 & 0.233 & 0.000 & 0.497 & 0.000 & 0.000 & 0.177 & 0.010 \\
LR & no\_peakflow & (3,7,7) & 355 & 0.005 & 0.109 & 0.000 & 0.806 & 0.014 & 0.200 & 0.400 & 0.032 \\
LR & no\_smartinhaler & (3,7,7) & 400 & 0.006 & 0.233 & 0.000 & 0.858 & 0.000 & 0.000 & 0.325 & 0.012 \\
LR & no\_smartwatch & (3,7,7) & 270 & 0.006 & 0.178 & 0.000 & 0.809 & 0.000 & 0.000 & 0.325 & 0.012 \\
LR & peakflow\_only & (3,7,7) & 131 & 0.006 & 0.223 & 0.000 & 0.551 & 0.000 & 0.000 & 0.149 & 0.009 \\
LR & smartinhaler\_only & (3,7,7) & 86 & 0.006 & 0.168 & 0.000 & 0.989 & 1.000 & 0.200 & 0.277 & 0.208 \\
LR & smartwatch\_only & (3,7,7) & 216 & 0.008 & 0.366 & 0.000 & 0.656 & 0.000 & 0.000 & 0.196 & 0.010 \\
LR & all & (3,7,14) & 427 & 0.005 & 0.050 & 0.000 & 0.475 & 0.005 & 0.200 & 0.310 & 0.011 \\
LR & environment\_only & (3,7,14) & 171 & 0.005 & 0.015 & 0.000 & 0.658 & 0.008 & 0.200 & 0.396 & 0.013 \\
LR & no\_environment & (3,7,14) & 315 & 0.006 & 0.213 & 0.000 & 0.458 & 0.000 & 0.000 & 0.162 & 0.010 \\
LR & no\_peakflow & (3,7,14) & 355 & 0.005 & 0.064 & 0.000 & 0.717 & 0.010 & 0.200 & 0.594 & 0.023 \\
LR & no\_smartinhaler & (3,7,14) & 400 & 0.005 & 0.074 & 0.000 & 0.525 & 0.000 & 0.000 & 0.257 & 0.011 \\
LR & no\_smartwatch & (3,7,14) & 270 & 0.005 & 0.035 & 0.000 & 0.436 & 0.005 & 0.200 & 0.331 & 0.013 \\
LR & peakflow\_only & (3,7,14) & 131 & 0.006 & 0.257 & 0.000 & 0.422 & 0.000 & 0.000 & 0.153 & 0.010 \\
LR & smartinhaler\_only & (3,7,14) & 86 & 0.005 & 0.087 & 0.000 & 0.822 & 0.016 & 0.200 & 0.212 & 0.012 \\
LR & smartwatch\_only & (3,7,14) & 216 & 0.008 & 0.411 & 0.000 & 0.750 & 0.000 & 0.000 & 0.241 & 0.011 \\
LR & all & (3,14,0) & 427 & 0.008 & 0.301 & 0.000 & 0.857 & 0.000 & 0.000 & 0.537 & 0.015 \\
LR & environment\_only & (3,14,0) & 171 & 0.007 & 0.230 & 0.000 & 0.854 & 0.000 & 0.000 & 0.592 & 0.017 \\
LR & no\_environment & (3,14,0) & 315 & 0.007 & 0.262 & 0.000 & 0.500 & 0.000 & 0.000 & 0.267 & 0.009 \\
LR & no\_peakflow & (3,14,0) & 355 & 0.008 & 0.279 & 0.000 & 0.883 & 0.000 & 0.000 & 0.388 & 0.011 \\
LR & no\_smartinhaler & (3,14,0) & 400 & 0.008 & 0.290 & 0.000 & 0.806 & 0.000 & 0.000 & 0.522 & 0.014 \\
LR & no\_smartwatch & (3,14,0) & 270 & 0.010 & 0.432 & 0.000 & 0.811 & 0.016 & 0.250 & 0.694 & 0.265 \\
LR & peakflow\_only & (3,14,0) & 131 & 0.006 & 0.169 & 0.000 & 0.406 & 0.005 & 0.250 & 0.280 & 0.010 \\
LR & smartinhaler\_only & (3,14,0) & 86 & 0.036 & 0.891 & 0.069 & 0.863 & 0.022 & 0.250 & 0.596 & 0.018 \\
LR & smartwatch\_only & (3,14,0) & 216 & 0.009 & 0.377 & 0.000 & 0.760 & 0.000 & 0.000 & 0.184 & 0.008 \\
LR & all & (3,14,7) & 427 & 0.007 & 0.268 & 0.000 & 0.651 & 0.000 & 0.000 & 0.416 & 0.012 \\
LR & environment\_only & (3,14,7) & 171 & 0.007 & 0.235 & 0.000 & 0.831 & 0.000 & 0.000 & 0.624 & 0.019 \\
LR & no\_environment & (3,14,7) & 315 & 0.008 & 0.311 & 0.000 & 0.509 & 0.000 & 0.000 & 0.240 & 0.009 \\
LR & no\_peakflow & (3,14,7) & 355 & 0.007 & 0.208 & 0.000 & 0.675 & 0.000 & 0.000 & 0.332 & 0.010 \\
LR & no\_smartinhaler & (3,14,7) & 400 & 0.007 & 0.213 & 0.000 & 0.592 & 0.000 & 0.000 & 0.415 & 0.012 \\
LR & no\_smartwatch & (3,14,7) & 270 & 0.010 & 0.437 & 0.000 & 0.725 & 0.011 & 0.250 & 0.584 & 0.025 \\
LR & peakflow\_only & (3,14,7) & 131 & 0.007 & 0.230 & 0.000 & 0.408 & 0.005 & 0.250 & 0.228 & 0.009 \\
LR & smartinhaler\_only & (3,14,7) & 86 & 0.023 & 0.803 & 0.000 & 0.763 & 0.000 & 0.000 & 0.412 & 0.012 \\
LR & smartwatch\_only & (3,14,7) & 216 & 0.007 & 0.257 & 0.000 & 0.624 & 0.000 & 0.000 & 0.251 & 0.009 \\
LR & all & (3,14,14) & 427 & 0.008 & 0.290 & 0.015 & 0.596 & 0.000 & 0.000 & 0.363 & 0.012 \\
LR & environment\_only & (3,14,14) & 171 & 0.009 & 0.383 & 0.000 & 0.790 & 0.000 & 0.000 & 0.532 & 0.016 \\
LR & no\_environment & (3,14,14) & 315 & 0.016 & 0.672 & 0.000 & 0.633 & 0.000 & 0.000 & 0.191 & 0.009 \\
LR & no\_peakflow & (3,14,14) & 355 & 0.007 & 0.262 & 0.000 & 0.731 & 0.000 & 0.000 & 0.427 & 0.013 \\
LR & no\_smartinhaler & (3,14,14) & 400 & 0.008 & 0.284 & 0.015 & 0.562 & 0.000 & 0.000 & 0.332 & 0.011 \\
LR & no\_smartwatch & (3,14,14) & 270 & 0.011 & 0.486 & 0.018 & 0.642 & 0.009 & 0.250 & 0.405 & 0.016 \\
LR & peakflow\_only & (3,14,14) & 131 & 0.010 & 0.492 & 0.000 & 0.404 & 0.005 & 0.250 & 0.141 & 0.009 \\
LR & smartinhaler\_only & (3,14,14) & 86 & 0.023 & 0.809 & 0.044 & 0.778 & 0.014 & 0.250 & 0.654 & 0.022 \\
LR & smartwatch\_only & (3,14,14) & 216 & 0.009 & 0.404 & 0.000 & 0.892 & 0.030 & 0.250 & 0.599 & 0.025 \\
LR & all & (4,3,0) & 427 & 0.008 & 0.407 & 0.000 & 0.885 & 0.023 & 0.200 & 0.393 & 0.019 \\
LR & environment\_only & (4,3,0) & 171 & 0.005 & 0.139 & 0.000 & 0.954 & 0.000 & 0.000 & 0.207 & 0.010 \\
LR & no\_environment & (4,3,0) & 315 & 0.006 & 0.245 & 0.000 & 0.810 & 0.013 & 0.200 & 0.282 & 0.011 \\
LR & no\_peakflow & (4,3,0) & 355 & 0.006 & 0.273 & 0.000 & 0.922 & 0.034 & 0.200 & 0.401 & 0.019 \\
LR & no\_smartinhaler & (4,3,0) & 400 & 0.008 & 0.417 & 0.000 & 0.910 & 0.000 & 0.000 & 0.368 & 0.015 \\
LR & no\_smartwatch & (4,3,0) & 270 & 0.006 & 0.222 & 0.000 & 0.949 & 0.000 & 0.000 & 0.368 & 0.014 \\
LR & peakflow\_only & (4,3,0) & 131 & 0.009 & 0.477 & 0.000 & 0.888 & 0.000 & 0.000 & 0.176 & 0.009 \\
LR & smartinhaler\_only & (4,3,0) & 86 & 0.006 & 0.222 & 0.000 & 0.988 & 0.000 & 0.000 & 0.122 & 0.008 \\
LR & smartwatch\_only & (4,3,0) & 216 & 0.007 & 0.375 & 0.000 & 0.871 & 0.020 & 0.200 & 0.319 & 0.019 \\
LR & all & (4,3,7) & 427 & 0.007 & 0.356 & 0.000 & 0.866 & 0.000 & 0.000 & 0.261 & 0.010 \\
LR & environment\_only & (4,3,7) & 171 & 0.005 & 0.079 & 0.000 & 0.975 & 0.000 & 0.000 & 0.195 & 0.009 \\
LR & no\_environment & (4,3,7) & 315 & 0.006 & 0.282 & 0.000 & 0.713 & 0.000 & 0.000 & 0.184 & 0.009 \\
LR & no\_peakflow & (4,3,7) & 355 & 0.006 & 0.222 & 0.000 & 0.907 & 0.000 & 0.000 & 0.220 & 0.009 \\
LR & no\_smartinhaler & (4,3,7) & 400 & 0.007 & 0.352 & 0.000 & 0.864 & 0.000 & 0.000 & 0.291 & 0.010 \\
LR & no\_smartwatch & (4,3,7) & 270 & 0.006 & 0.181 & 0.000 & 0.955 & 0.000 & 0.000 & 0.345 & 0.011 \\
LR & peakflow\_only & (4,3,7) & 131 & 0.005 & 0.123 & 0.000 & 0.912 & 0.000 & 0.000 & 0.157 & 0.009 \\
LR & smartinhaler\_only & (4,3,7) & 86 & 0.005 & 0.120 & 0.000 & 0.987 & 0.000 & 0.000 & 0.200 & 0.009 \\
LR & smartwatch\_only & (4,3,7) & 216 & 0.008 & 0.472 & 0.000 & 0.889 & 0.000 & 0.000 & 0.217 & 0.009 \\
LR & all & (4,3,14) & 427 & 0.006 & 0.269 & 0.000 & 0.712 & 0.000 & 0.000 & 0.266 & 0.010 \\
LR & environment\_only & (4,3,14) & 171 & 0.005 & 0.074 & 0.000 & 0.881 & 0.000 & 0.000 & 0.185 & 0.009 \\
LR & no\_environment & (4,3,14) & 315 & 0.006 & 0.204 & 0.000 & 0.590 & 0.000 & 0.000 & 0.170 & 0.009 \\
LR & no\_peakflow & (4,3,14) & 355 & 0.006 & 0.222 & 0.000 & 0.790 & 0.000 & 0.000 & 0.260 & 0.010 \\
LR & no\_smartinhaler & (4,3,14) & 400 & 0.006 & 0.259 & 0.000 & 0.722 & 0.000 & 0.000 & 0.255 & 0.010 \\
LR & no\_smartwatch & (4,3,14) & 270 & 0.006 & 0.199 & 0.000 & 0.805 & 0.000 & 0.000 & 0.291 & 0.011 \\
LR & peakflow\_only & (4,3,14) & 131 & 0.005 & 0.090 & 0.000 & 0.481 & 0.000 & 0.000 & 0.152 & 0.009 \\
LR & smartinhaler\_only & (4,3,14) & 86 & 0.005 & 0.116 & 0.000 & 0.987 & 0.000 & 0.000 & 0.138 & 0.009 \\
LR & smartwatch\_only & (4,3,14) & 216 & 0.007 & 0.384 & 0.000 & 0.855 & 0.000 & 0.000 & 0.297 & 0.011 \\
LR & all & (4,7,0) & 427 & 0.006 & 0.233 & 0.000 & 0.798 & 0.000 & 0.000 & 0.235 & 0.010 \\
LR & environment\_only & (4,7,0) & 171 & 0.005 & 0.020 & 0.000 & 0.788 & 0.013 & 0.200 & 0.338 & 0.012 \\
LR & no\_environment & (4,7,0) & 315 & 0.007 & 0.302 & 0.000 & 0.557 & 0.000 & 0.000 & 0.193 & 0.009 \\
LR & no\_peakflow & (4,7,0) & 355 & 0.006 & 0.114 & 0.000 & 0.832 & 0.000 & 0.000 & 0.216 & 0.009 \\
LR & no\_smartinhaler & (4,7,0) & 400 & 0.006 & 0.238 & 0.000 & 0.819 & 0.000 & 0.000 & 0.282 & 0.010 \\
LR & no\_smartwatch & (4,7,0) & 270 & 0.006 & 0.119 & 0.000 & 0.870 & 0.000 & 0.000 & 0.306 & 0.011 \\
LR & peakflow\_only & (4,7,0) & 131 & 0.007 & 0.292 & 0.000 & 0.858 & 0.000 & 0.000 & 0.163 & 0.009 \\
LR & smartinhaler\_only & (4,7,0) & 86 & 0.006 & 0.163 & 0.000 & 0.953 & 0.000 & 0.000 & 0.189 & 0.009 \\
LR & smartwatch\_only & (4,7,0) & 216 & 0.012 & 0.584 & 0.000 & 0.769 & 0.000 & 0.000 & 0.183 & 0.009 \\
LR & all & (4,7,7) & 427 & 0.007 & 0.257 & 0.000 & 0.826 & 0.000 & 0.000 & 0.416 & 0.013 \\
LR & environment\_only & (4,7,7) & 171 & 0.005 & 0.015 & 0.000 & 0.697 & 0.009 & 0.200 & 0.301 & 0.012 \\
LR & no\_environment & (4,7,7) & 315 & 0.008 & 0.371 & 0.000 & 0.493 & 0.000 & 0.000 & 0.191 & 0.010 \\
LR & no\_peakflow & (4,7,7) & 355 & 0.005 & 0.084 & 0.000 & 0.815 & 0.015 & 0.200 & 0.416 & 0.016 \\
LR & no\_smartinhaler & (4,7,7) & 400 & 0.006 & 0.248 & 0.000 & 0.847 & 0.000 & 0.000 & 0.422 & 0.013 \\
LR & no\_smartwatch & (4,7,7) & 270 & 0.006 & 0.144 & 0.000 & 0.839 & 0.000 & 0.000 & 0.302 & 0.011 \\
LR & peakflow\_only & (4,7,7) & 131 & 0.009 & 0.441 & 0.000 & 0.692 & 0.000 & 0.000 & 0.123 & 0.009 \\
LR & smartinhaler\_only & (4,7,7) & 86 & 0.006 & 0.196 & 0.000 & 0.850 & 0.000 & 0.000 & 0.196 & 0.010 \\
LR & smartwatch\_only & (4,7,7) & 216 & 0.010 & 0.490 & 0.000 & 0.684 & 0.000 & 0.000 & 0.258 & 0.011 \\
LR & all & (4,7,14) & 427 & 0.005 & 0.035 & 0.000 & 0.418 & 0.000 & 0.000 & 0.304 & 0.011 \\
LR & environment\_only & (4,7,14) & 171 & 0.005 & 0.005 & 0.000 & 0.765 & 0.000 & 0.000 & 0.469 & 0.015 \\
LR & no\_environment & (4,7,14) & 315 & 0.006 & 0.188 & 0.000 & 0.457 & 0.000 & 0.000 & 0.131 & 0.009 \\
LR & no\_peakflow & (4,7,14) & 355 & 0.005 & 0.030 & 0.000 & 0.787 & 0.026 & 0.400 & 0.624 & 0.022 \\
LR & no\_smartinhaler & (4,7,14) & 400 & 0.005 & 0.035 & 0.000 & 0.468 & 0.000 & 0.000 & 0.260 & 0.011 \\
LR & no\_smartwatch & (4,7,14) & 270 & 0.005 & 0.010 & 0.000 & 0.465 & 0.005 & 0.200 & 0.344 & 0.014 \\
LR & peakflow\_only & (4,7,14) & 131 & 0.005 & 0.168 & 0.000 & 0.435 & 0.000 & 0.000 & 0.106 & 0.009 \\
LR & smartinhaler\_only & (4,7,14) & 86 & 0.005 & 0.052 & 0.000 & 0.856 & 0.000 & 0.000 & 0.245 & 0.011 \\
LR & smartwatch\_only & (4,7,14) & 216 & 0.006 & 0.238 & 0.000 & 0.853 & 0.000 & 0.000 & 0.465 & 0.015 \\
LR & all & (4,14,0) & 427 & 0.007 & 0.273 & 0.000 & 0.911 & 0.034 & 0.250 & 0.692 & 0.028 \\
LR & environment\_only & (4,14,0) & 171 & 0.007 & 0.197 & 0.000 & 0.866 & 0.000 & 0.000 & 0.704 & 0.024 \\
LR & no\_environment & (4,14,0) & 315 & 0.008 & 0.279 & 0.000 & 0.657 & 0.008 & 0.250 & 0.334 & 0.010 \\
LR & no\_peakflow & (4,14,0) & 355 & 0.007 & 0.202 & 0.000 & 0.929 & 0.000 & 0.000 & 0.496 & 0.015 \\
LR & no\_smartinhaler & (4,14,0) & 400 & 0.007 & 0.230 & 0.000 & 0.820 & 0.016 & 0.250 & 0.608 & 0.022 \\
LR & no\_smartwatch & (4,14,0) & 270 & 0.010 & 0.470 & 0.000 & 0.883 & 0.026 & 0.250 & 0.840 & 0.279 \\
LR & peakflow\_only & (4,14,0) & 131 & 0.007 & 0.262 & 0.000 & 0.583 & 0.000 & 0.000 & 0.243 & 0.009 \\
LR & smartinhaler\_only & (4,14,0) & 86 & 0.036 & 0.891 & 0.069 & 0.834 & 0.018 & 0.250 & 0.671 & 0.024 \\
LR & smartwatch\_only & (4,14,0) & 216 & 0.010 & 0.464 & 0.000 & 0.834 & 0.000 & 0.000 & 0.160 & 0.008 \\
LR & all & (4,14,7) & 427 & 0.009 & 0.410 & 0.000 & 0.838 & 0.019 & 0.250 & 0.648 & 0.021 \\
LR & environment\_only & (4,14,7) & 171 & 0.008 & 0.344 & 0.000 & 0.841 & 0.000 & 0.000 & 0.719 & 0.025 \\
LR & no\_environment & (4,14,7) & 315 & 0.007 & 0.246 & 0.000 & 0.493 & 0.000 & 0.000 & 0.229 & 0.009 \\
LR & no\_peakflow & (4,14,7) & 355 & 0.007 & 0.251 & 0.000 & 0.876 & 0.000 & 0.000 & 0.549 & 0.017 \\
LR & no\_smartinhaler & (4,14,7) & 400 & 0.008 & 0.333 & 0.000 & 0.746 & 0.012 & 0.250 & 0.560 & 0.018 \\
LR & no\_smartwatch & (4,14,7) & 270 & 0.012 & 0.546 & 0.000 & 0.808 & 0.016 & 0.250 & 0.784 & 0.033 \\
LR & peakflow\_only & (4,14,7) & 131 & 0.008 & 0.328 & 0.000 & 0.572 & 0.000 & 0.000 & 0.275 & 0.010 \\
LR & smartinhaler\_only & (4,14,7) & 86 & 0.023 & 0.803 & 0.044 & 0.717 & 0.000 & 0.000 & 0.571 & 0.017 \\
LR & smartwatch\_only & (4,14,7) & 216 & 0.010 & 0.437 & 0.000 & 0.832 & 0.000 & 0.000 & 0.171 & 0.009 \\
LR & all & (4,14,14) & 427 & 0.010 & 0.454 & 0.000 & 0.822 & 0.018 & 0.250 & 0.607 & 0.019 \\
LR & environment\_only & (4,14,14) & 171 & 0.008 & 0.355 & 0.000 & 0.818 & 0.000 & 0.000 & 0.572 & 0.017 \\
LR & no\_environment & (4,14,14) & 315 & 0.019 & 0.721 & 0.000 & 0.843 & 0.000 & 0.000 & 0.245 & 0.010 \\
LR & no\_peakflow & (4,14,14) & 355 & 0.007 & 0.208 & 0.000 & 0.803 & 0.000 & 0.000 & 0.460 & 0.015 \\
LR & no\_smartinhaler & (4,14,14) & 400 & 0.010 & 0.432 & 0.000 & 0.818 & 0.018 & 0.250 & 0.586 & 0.021 \\
LR & no\_smartwatch & (4,14,14) & 270 & 0.013 & 0.579 & 0.025 & 0.809 & 0.017 & 0.250 & 0.732 & 0.027 \\
LR & peakflow\_only & (4,14,14) & 131 & 0.023 & 0.787 & 0.000 & 0.812 & 0.000 & 0.000 & 0.211 & 0.009 \\
LR & smartinhaler\_only & (4,14,14) & 86 & 0.023 & 0.803 & 0.000 & 0.812 & 0.000 & 0.000 & 0.702 & 0.025 \\
LR & smartwatch\_only & (4,14,14) & 216 & 0.008 & 0.317 & 0.000 & 0.889 & 0.000 & 0.000 & 0.528 & 0.022 \\
RF & all & (2,3,0) & 427 & 0.016 & 0.727 & 0.000 & 0.993 & 0.000 & 0.000 & 0.516 & 0.012 \\
RF & environment\_only & (2,3,0) & 171 & 0.011 & 0.606 & 0.000 & 0.993 & 0.000 & 0.000 & 0.569 & 0.016 \\
RF & no\_environment & (2,3,0) & 315 & 0.013 & 0.671 & 0.000 & 0.993 & 0.000 & 0.000 & 0.156 & 0.006 \\
RF & no\_peakflow & (2,3,0) & 355 & 0.013 & 0.671 & 0.000 & 0.993 & 0.000 & 0.000 & 0.508 & 0.012 \\
RF & no\_smartinhaler & (2,3,0) & 400 & 0.020 & 0.796 & 0.000 & 0.993 & 0.000 & 0.000 & 0.398 & 0.008 \\
RF & no\_smartwatch & (2,3,0) & 270 & 0.008 & 0.472 & 0.000 & 0.993 & 0.000 & 0.000 & 0.524 & 0.022 \\
RF & peakflow\_only & (2,3,0) & 131 & 0.011 & 0.669 & 0.000 & 0.993 & 0.000 & 0.000 & 0.342 & 0.007 \\
RF & smartinhaler\_only & (2,3,0) & 86 & 0.006 & 0.287 & 0.000 & 0.993 & 0.000 & 0.000 & 0.329 & 0.007 \\
RF & smartwatch\_only & (2,3,0) & 216 & 0.013 & 0.681 & 0.000 & 0.993 & 0.000 & 0.000 & 0.550 & 0.012 \\
RF & all & (2,3,7) & 427 & 0.015 & 0.708 & 0.000 & 0.992 & 0.000 & 0.000 & 0.395 & 0.008 \\
RF & environment\_only & (2,3,7) & 171 & 0.008 & 0.431 & 0.000 & 0.992 & 0.000 & 0.000 & 0.607 & 0.022 \\
RF & no\_environment & (2,3,7) & 315 & 0.009 & 0.486 & 0.000 & 0.992 & 0.000 & 0.000 & 0.286 & 0.007 \\
RF & no\_peakflow & (2,3,7) & 355 & 0.010 & 0.551 & 0.000 & 0.992 & 0.000 & 0.000 & 0.610 & 0.045 \\
RF & no\_smartinhaler & (2,3,7) & 400 & 0.010 & 0.565 & 0.000 & 0.992 & 0.000 & 0.000 & 0.483 & 0.010 \\
RF & no\_smartwatch & (2,3,7) & 270 & 0.007 & 0.366 & 0.000 & 0.992 & 0.000 & 0.000 & 0.532 & 0.020 \\
RF & peakflow\_only & (2,3,7) & 131 & 0.008 & 0.475 & 0.000 & 0.992 & 0.000 & 0.000 & 0.344 & 0.008 \\
RF & smartinhaler\_only & (2,3,7) & 86 & 0.006 & 0.282 & 0.000 & 0.992 & 0.000 & 0.000 & 0.559 & 0.011 \\
RF & smartwatch\_only & (2,3,7) & 216 & 0.011 & 0.588 & 0.000 & 0.992 & 0.000 & 0.000 & 0.667 & 0.031 \\
RF & all & (2,3,14) & 427 & 0.008 & 0.440 & 0.000 & 0.992 & 0.000 & 0.000 & 0.546 & 0.013 \\
RF & environment\_only & (2,3,14) & 171 & 0.007 & 0.343 & 0.000 & 0.992 & 0.000 & 0.000 & 0.573 & 0.022 \\
RF & no\_environment & (2,3,14) & 315 & 0.006 & 0.292 & 0.000 & 0.992 & 0.000 & 0.000 & 0.310 & 0.007 \\
RF & no\_peakflow & (2,3,14) & 355 & 0.006 & 0.241 & 0.000 & 0.992 & 0.000 & 0.000 & 0.675 & 0.039 \\
RF & no\_smartinhaler & (2,3,14) & 400 & 0.009 & 0.491 & 0.000 & 0.992 & 0.000 & 0.000 & 0.456 & 0.010 \\
RF & no\_smartwatch & (2,3,14) & 270 & 0.006 & 0.259 & 0.000 & 0.990 & 0.000 & 0.000 & 0.483 & 0.013 \\
RF & peakflow\_only & (2,3,14) & 131 & 0.006 & 0.257 & 0.000 & 0.990 & 0.000 & 0.000 & 0.243 & 0.007 \\
RF & smartinhaler\_only & (2,3,14) & 86 & 0.006 & 0.282 & 0.000 & 0.992 & 0.000 & 0.000 & 0.590 & 0.012 \\
RF & smartwatch\_only & (2,3,14) & 216 & 0.006 & 0.255 & 0.000 & 0.992 & 0.000 & 0.000 & 0.688 & 0.044 \\
RF & all & (2,7,0) & 427 & 0.012 & 0.594 & 0.000 & 0.992 & 0.000 & 0.000 & 0.525 & 0.054 \\
RF & environment\_only & (2,7,0) & 171 & 0.010 & 0.535 & 0.000 & 0.992 & 0.000 & 0.000 & 0.590 & 0.074 \\
RF & no\_environment & (2,7,0) & 315 & 0.007 & 0.307 & 0.000 & 0.992 & 0.000 & 0.000 & 0.460 & 0.010 \\
RF & no\_peakflow & (2,7,0) & 355 & 0.012 & 0.599 & 0.000 & 0.992 & 0.000 & 0.000 & 0.584 & 0.118 \\
RF & no\_smartinhaler & (2,7,0) & 400 & 0.011 & 0.574 & 0.000 & 0.992 & 0.000 & 0.000 & 0.474 & 0.036 \\
RF & no\_smartwatch & (2,7,0) & 270 & 0.009 & 0.441 & 0.000 & 0.992 & 0.000 & 0.000 & 0.512 & 0.015 \\
RF & peakflow\_only & (2,7,0) & 131 & 0.006 & 0.297 & 0.000 & 0.992 & 0.000 & 0.000 & 0.320 & 0.009 \\
RF & smartinhaler\_only & (2,7,0) & 86 & 0.017 & 0.760 & 0.000 & 0.992 & 0.000 & 0.000 & 0.619 & 0.014 \\
RF & smartwatch\_only & (2,7,0) & 216 & 0.009 & 0.495 & 0.000 & 0.992 & 0.000 & 0.000 & 0.714 & 0.149 \\
RF & all & (2,7,7) & 427 & 0.010 & 0.520 & 0.000 & 0.992 & 0.000 & 0.000 & 0.476 & 0.020 \\
RF & environment\_only & (2,7,7) & 171 & 0.007 & 0.302 & 0.000 & 0.992 & 0.000 & 0.000 & 0.635 & 0.076 \\
RF & no\_environment & (2,7,7) & 315 & 0.007 & 0.262 & 0.000 & 0.992 & 0.000 & 0.000 & 0.373 & 0.009 \\
RF & no\_peakflow & (2,7,7) & 355 & 0.011 & 0.550 & 0.000 & 0.992 & 0.000 & 0.000 & 0.649 & 0.120 \\
RF & no\_smartinhaler & (2,7,7) & 400 & 0.010 & 0.530 & 0.000 & 0.992 & 0.000 & 0.000 & 0.490 & 0.026 \\
RF & no\_smartwatch & (2,7,7) & 270 & 0.007 & 0.307 & 0.000 & 0.992 & 0.000 & 0.000 & 0.465 & 0.010 \\
RF & peakflow\_only & (2,7,7) & 131 & 0.007 & 0.396 & 0.000 & 0.992 & 0.000 & 0.000 & 0.281 & 0.008 \\
RF & smartinhaler\_only & (2,7,7) & 86 & 0.018 & 0.780 & 0.000 & 0.992 & 0.000 & 0.000 & 0.621 & 0.014 \\
RF & smartwatch\_only & (2,7,7) & 216 & 0.008 & 0.391 & 0.000 & 0.992 & 0.000 & 0.000 & 0.698 & 0.132 \\
RF & all & (2,7,14) & 427 & 0.009 & 0.465 & 0.000 & 0.992 & 0.000 & 0.000 & 0.554 & 0.090 \\
RF & environment\_only & (2,7,14) & 171 & 0.006 & 0.208 & 0.000 & 0.992 & 0.000 & 0.000 & 0.540 & 0.019 \\
RF & no\_environment & (2,7,14) & 315 & 0.006 & 0.203 & 0.000 & 0.980 & 0.000 & 0.000 & 0.494 & 0.012 \\
RF & no\_peakflow & (2,7,14) & 355 & 0.007 & 0.312 & 0.000 & 0.992 & 0.000 & 0.000 & 0.685 & 0.346 \\
RF & no\_smartinhaler & (2,7,14) & 400 & 0.008 & 0.396 & 0.000 & 0.994 & 1.000 & 0.333 & 0.574 & 0.341 \\
RF & no\_smartwatch & (2,7,14) & 270 & 0.007 & 0.322 & 0.000 & 0.958 & 0.000 & 0.000 & 0.433 & 0.010 \\
RF & peakflow\_only & (2,7,14) & 131 & 0.006 & 0.267 & 0.000 & 0.634 & 0.000 & 0.000 & 0.271 & 0.008 \\
RF & smartinhaler\_only & (2,7,14) & 86 & 0.005 & 0.146 & 0.000 & 0.992 & 0.000 & 0.000 & 0.617 & 0.016 \\
RF & smartwatch\_only & (2,7,14) & 216 & 0.006 & 0.248 & 0.000 & 0.992 & 0.000 & 0.000 & 0.705 & 0.159 \\
RF & all & (2,14,0) & 427 & 0.200 & 0.978 & 0.000 & 0.994 & 0.000 & 0.000 & 0.545 & 0.009 \\
RF & environment\_only & (2,14,0) & 171 & 0.011 & 0.525 & 0.000 & 0.994 & 0.000 & 0.000 & 0.488 & 0.008 \\
RF & no\_environment & (2,14,0) & 315 & 0.007 & 0.186 & 0.000 & 0.994 & 0.000 & 0.000 & 0.428 & 0.007 \\
RF & no\_peakflow & (2,14,0) & 355 & 0.036 & 0.852 & 0.000 & 0.994 & 0.000 & 0.000 & 0.556 & 0.010 \\
RF & no\_smartinhaler & (2,14,0) & 400 & 0.030 & 0.825 & 0.000 & 0.994 & 0.000 & 0.000 & 0.448 & 0.008 \\
RF & no\_smartwatch & (2,14,0) & 270 & 0.008 & 0.344 & 0.000 & 0.986 & 0.000 & 0.000 & 0.493 & 0.011 \\
RF & peakflow\_only & (2,14,0) & 131 & 0.006 & 0.098 & 0.000 & 0.721 & 0.010 & 0.500 & 0.493 & 0.032 \\
RF & smartinhaler\_only & (2,14,0) & 86 & 0.006 & 0.142 & 0.000 & 0.891 & 0.000 & 0.000 & 0.665 & 0.013 \\
RF & smartwatch\_only & (2,14,0) & 216 & 0.006 & 0.158 & 0.000 & 0.994 & 0.000 & 0.000 & 0.574 & 0.014 \\
RF & all & (2,14,7) & 427 & 0.024 & 0.781 & 0.000 & 0.994 & 0.000 & 0.000 & 0.466 & 0.008 \\
RF & environment\_only & (2,14,7) & 171 & 0.007 & 0.208 & 0.000 & 0.994 & 0.000 & 0.000 & 0.452 & 0.008 \\
RF & no\_environment & (2,14,7) & 315 & 0.006 & 0.066 & 0.000 & 0.994 & 0.000 & 0.000 & 0.114 & 0.005 \\
RF & no\_peakflow & (2,14,7) & 355 & 0.017 & 0.683 & 0.000 & 0.994 & 0.000 & 0.000 & 0.552 & 0.010 \\
RF & no\_smartinhaler & (2,14,7) & 400 & 0.021 & 0.743 & 0.000 & 0.994 & 0.000 & 0.000 & 0.434 & 0.008 \\
RF & no\_smartwatch & (2,14,7) & 270 & 0.007 & 0.257 & 0.000 & 0.994 & 0.000 & 0.000 & 0.546 & 0.013 \\
RF & peakflow\_only & (2,14,7) & 131 & 0.008 & 0.311 & 0.000 & 0.690 & 0.010 & 0.500 & 0.476 & 0.026 \\
RF & smartinhaler\_only & (2,14,7) & 86 & 0.026 & 0.836 & 0.000 & 0.994 & 0.000 & 0.000 & 0.672 & 0.013 \\
RF & smartwatch\_only & (2,14,7) & 216 & 0.006 & 0.087 & 0.000 & 0.994 & 0.000 & 0.000 & 0.455 & 0.012 \\
RF & all & (2,14,14) & 427 & 0.017 & 0.683 & 0.000 & 0.994 & 0.000 & 0.000 & 0.613 & 0.012 \\
RF & environment\_only & (2,14,14) & 171 & 0.009 & 0.421 & 0.000 & 0.994 & 0.000 & 0.000 & 0.550 & 0.010 \\
RF & no\_environment & (2,14,14) & 315 & 0.006 & 0.120 & 0.000 & 0.988 & 0.000 & 0.000 & 0.641 & 0.015 \\
RF & no\_peakflow & (2,14,14) & 355 & 0.018 & 0.705 & 0.000 & 0.994 & 0.000 & 0.000 & 0.672 & 0.255 \\
RF & no\_smartinhaler & (2,14,14) & 400 & 0.013 & 0.596 & 0.000 & 0.994 & 0.000 & 0.000 & 0.789 & 0.021 \\
RF & no\_smartwatch & (2,14,14) & 270 & 0.007 & 0.240 & 0.000 & 0.829 & 0.018 & 0.500 & 0.594 & 0.017 \\
RF & peakflow\_only & (2,14,14) & 131 & 0.005 & 0.022 & 0.000 & 0.394 & 0.005 & 0.500 & 0.523 & 0.031 \\
RF & smartinhaler\_only & (2,14,14) & 86 & 0.010 & 0.470 & 0.000 & 0.866 & 0.023 & 0.500 & 0.764 & 0.022 \\
RF & smartwatch\_only & (2,14,14) & 216 & 0.006 & 0.093 & 0.000 & 0.994 & 0.000 & 0.000 & 0.589 & 0.504 \\
RF & all & (3,3,0) & 427 & 0.016 & 0.731 & 0.000 & 0.988 & 0.000 & 0.000 & 0.457 & 0.015 \\
RF & environment\_only & (3,3,0) & 171 & 0.015 & 0.713 & 0.000 & 0.988 & 0.000 & 0.000 & 0.464 & 0.019 \\
RF & no\_environment & (3,3,0) & 315 & 0.009 & 0.509 & 0.000 & 0.988 & 0.000 & 0.000 & 0.378 & 0.011 \\
RF & no\_peakflow & (3,3,0) & 355 & 0.011 & 0.606 & 0.000 & 0.988 & 0.000 & 0.000 & 0.504 & 0.018 \\
RF & no\_smartinhaler & (3,3,0) & 400 & 0.019 & 0.773 & 0.000 & 0.988 & 0.000 & 0.000 & 0.401 & 0.012 \\
RF & no\_smartwatch & (3,3,0) & 270 & 0.015 & 0.713 & 0.000 & 0.988 & 0.000 & 0.000 & 0.414 & 0.016 \\
RF & peakflow\_only & (3,3,0) & 131 & 0.013 & 0.738 & 0.000 & 0.988 & 0.000 & 0.000 & 0.305 & 0.010 \\
RF & smartinhaler\_only & (3,3,0) & 86 & 0.006 & 0.282 & 0.000 & 0.988 & 0.000 & 0.000 & 0.343 & 0.011 \\
RF & smartwatch\_only & (3,3,0) & 216 & 0.009 & 0.532 & 0.000 & 0.988 & 0.000 & 0.000 & 0.599 & 0.020 \\
RF & all & (3,3,7) & 427 & 0.009 & 0.519 & 0.000 & 0.987 & 0.000 & 0.000 & 0.414 & 0.013 \\
RF & environment\_only & (3,3,7) & 171 & 0.007 & 0.384 & 0.000 & 0.987 & 0.000 & 0.000 & 0.420 & 0.017 \\
RF & no\_environment & (3,3,7) & 315 & 0.008 & 0.458 & 0.000 & 0.987 & 0.000 & 0.000 & 0.386 & 0.012 \\
RF & no\_peakflow & (3,3,7) & 355 & 0.009 & 0.486 & 0.000 & 0.987 & 0.000 & 0.000 & 0.505 & 0.034 \\
RF & no\_smartinhaler & (3,3,7) & 400 & 0.009 & 0.500 & 0.000 & 0.987 & 0.000 & 0.000 & 0.327 & 0.012 \\
RF & no\_smartwatch & (3,3,7) & 270 & 0.006 & 0.292 & 0.000 & 0.987 & 0.000 & 0.000 & 0.393 & 0.015 \\
RF & peakflow\_only & (3,3,7) & 131 & 0.006 & 0.373 & 0.000 & 0.987 & 0.000 & 0.000 & 0.293 & 0.010 \\
RF & smartinhaler\_only & (3,3,7) & 86 & 0.006 & 0.306 & 0.000 & 0.987 & 0.000 & 0.000 & 0.386 & 0.012 \\
RF & smartwatch\_only & (3,3,7) & 216 & 0.010 & 0.574 & 0.000 & 0.987 & 0.000 & 0.000 & 0.411 & 0.013 \\
RF & all & (3,3,14) & 427 & 0.007 & 0.380 & 0.000 & 0.987 & 0.000 & 0.000 & 0.439 & 0.015 \\
RF & environment\_only & (3,3,14) & 171 & 0.007 & 0.329 & 0.000 & 0.987 & 0.000 & 0.000 & 0.480 & 0.020 \\
RF & no\_environment & (3,3,14) & 315 & 0.006 & 0.204 & 0.000 & 0.987 & 0.000 & 0.000 & 0.423 & 0.013 \\
RF & no\_peakflow & (3,3,14) & 355 & 0.007 & 0.356 & 0.000 & 0.987 & 0.000 & 0.000 & 0.487 & 0.020 \\
RF & no\_smartinhaler & (3,3,14) & 400 & 0.007 & 0.366 & 0.000 & 0.987 & 0.000 & 0.000 & 0.512 & 0.023 \\
RF & no\_smartwatch & (3,3,14) & 270 & 0.006 & 0.292 & 0.000 & 0.984 & 0.000 & 0.000 & 0.413 & 0.027 \\
RF & peakflow\_only & (3,3,14) & 131 & 0.006 & 0.340 & 0.000 & 0.987 & 0.000 & 0.000 & 0.342 & 0.011 \\
RF & smartinhaler\_only & (3,3,14) & 86 & 0.006 & 0.282 & 0.000 & 0.987 & 0.000 & 0.000 & 0.647 & 0.021 \\
RF & smartwatch\_only & (3,3,14) & 216 & 0.005 & 0.181 & 0.000 & 0.987 & 0.000 & 0.000 & 0.771 & 0.035 \\
RF & all & (3,7,0) & 427 & 0.018 & 0.738 & 0.000 & 0.987 & 0.000 & 0.000 & 0.551 & 0.080 \\
RF & environment\_only & (3,7,0) & 171 & 0.011 & 0.550 & 0.000 & 0.987 & 0.000 & 0.000 & 0.416 & 0.050 \\
RF & no\_environment & (3,7,0) & 315 & 0.007 & 0.297 & 0.000 & 0.987 & 0.000 & 0.000 & 0.492 & 0.015 \\
RF & no\_peakflow & (3,7,0) & 355 & 0.015 & 0.688 & 0.000 & 0.987 & 0.000 & 0.000 & 0.582 & 0.080 \\
RF & no\_smartinhaler & (3,7,0) & 400 & 0.019 & 0.743 & 0.000 & 0.987 & 0.000 & 0.000 & 0.522 & 0.078 \\
RF & no\_smartwatch & (3,7,0) & 270 & 0.010 & 0.505 & 0.000 & 0.987 & 0.000 & 0.000 & 0.404 & 0.032 \\
RF & peakflow\_only & (3,7,0) & 131 & 0.021 & 0.832 & 0.000 & 0.987 & 0.000 & 0.000 & 0.229 & 0.011 \\
RF & smartinhaler\_only & (3,7,0) & 86 & 0.016 & 0.745 & 0.000 & 0.987 & 0.000 & 0.000 & 0.610 & 0.020 \\
RF & smartwatch\_only & (3,7,0) & 216 & 0.008 & 0.411 & 0.000 & 0.987 & 0.000 & 0.000 & 0.750 & 0.103 \\
RF & all & (3,7,7) & 427 & 0.021 & 0.777 & 0.000 & 0.987 & 0.000 & 0.000 & 0.412 & 0.024 \\
RF & environment\_only & (3,7,7) & 171 & 0.006 & 0.248 & 0.000 & 0.987 & 0.000 & 0.000 & 0.404 & 0.025 \\
RF & no\_environment & (3,7,7) & 315 & 0.006 & 0.243 & 0.000 & 0.987 & 0.000 & 0.000 & 0.499 & 0.016 \\
RF & no\_peakflow & (3,7,7) & 355 & 0.011 & 0.545 & 0.000 & 0.987 & 0.000 & 0.000 & 0.571 & 0.080 \\
RF & no\_smartinhaler & (3,7,7) & 400 & 0.011 & 0.564 & 0.000 & 0.987 & 0.000 & 0.000 & 0.424 & 0.030 \\
RF & no\_smartwatch & (3,7,7) & 270 & 0.006 & 0.243 & 0.000 & 0.987 & 0.000 & 0.000 & 0.374 & 0.012 \\
RF & peakflow\_only & (3,7,7) & 131 & 0.016 & 0.762 & 0.000 & 0.987 & 0.000 & 0.000 & 0.237 & 0.011 \\
RF & smartinhaler\_only & (3,7,7) & 86 & 0.012 & 0.636 & 0.000 & 0.987 & 0.000 & 0.000 & 0.499 & 0.017 \\
RF & smartwatch\_only & (3,7,7) & 216 & 0.007 & 0.287 & 0.000 & 0.987 & 0.000 & 0.000 & 0.730 & 0.109 \\
RF & all & (3,7,14) & 427 & 0.020 & 0.757 & 0.000 & 0.981 & 0.000 & 0.000 & 0.565 & 0.081 \\
RF & environment\_only & (3,7,14) & 171 & 0.006 & 0.178 & 0.000 & 0.986 & 0.000 & 0.000 & 0.528 & 0.053 \\
RF & no\_environment & (3,7,14) & 315 & 0.008 & 0.396 & 0.000 & 0.986 & 0.000 & 0.000 & 0.438 & 0.015 \\
RF & no\_peakflow & (3,7,14) & 355 & 0.008 & 0.376 & 0.000 & 0.986 & 0.000 & 0.000 & 0.688 & 0.087 \\
RF & no\_smartinhaler & (3,7,14) & 400 & 0.013 & 0.629 & 0.000 & 0.981 & 0.000 & 0.000 & 0.412 & 0.060 \\
RF & no\_smartwatch & (3,7,14) & 270 & 0.006 & 0.233 & 0.000 & 0.986 & 0.000 & 0.000 & 0.317 & 0.012 \\
RF & peakflow\_only & (3,7,14) & 131 & 0.009 & 0.525 & 0.000 & 0.956 & 0.000 & 0.000 & 0.119 & 0.009 \\
RF & smartinhaler\_only & (3,7,14) & 86 & 0.013 & 0.661 & 0.000 & 0.986 & 0.000 & 0.000 & 0.465 & 0.015 \\
RF & smartwatch\_only & (3,7,14) & 216 & 0.006 & 0.248 & 0.000 & 0.986 & 0.000 & 0.000 & 0.772 & 0.105 \\
RF & all & (3,14,0) & 427 & 0.250 & 0.984 & 0.000 & 0.989 & 0.000 & 0.000 & 0.374 & 0.011 \\
RF & environment\_only & (3,14,0) & 171 & 0.009 & 0.383 & 0.000 & 0.989 & 0.000 & 0.000 & 0.414 & 0.012 \\
RF & no\_environment & (3,14,0) & 315 & 0.006 & 0.148 & 0.000 & 0.989 & 0.000 & 0.000 & 0.474 & 0.013 \\
RF & no\_peakflow & (3,14,0) & 355 & 0.125 & 0.962 & 0.000 & 0.989 & 0.000 & 0.000 & 0.648 & 0.021 \\
RF & no\_smartinhaler & (3,14,0) & 400 & 0.143 & 0.967 & 0.000 & 0.989 & 0.000 & 0.000 & 0.397 & 0.011 \\
RF & no\_smartwatch & (3,14,0) & 270 & 0.007 & 0.246 & 0.000 & 0.989 & 0.000 & 0.000 & 0.403 & 0.015 \\
RF & peakflow\_only & (3,14,0) & 131 & 0.008 & 0.311 & 0.000 & 0.834 & 0.018 & 0.250 & 0.314 & 0.028 \\
RF & smartinhaler\_only & (3,14,0) & 86 & 0.030 & 0.863 & 0.000 & 0.886 & 0.000 & 0.000 & 0.624 & 0.020 \\
RF & smartwatch\_only & (3,14,0) & 216 & 0.006 & 0.120 & 0.000 & 0.989 & 0.000 & 0.000 & 0.718 & 0.049 \\
RF & all & (3,14,7) & 427 & 0.056 & 0.907 & 0.000 & 0.988 & 0.000 & 0.000 & 0.498 & 0.014 \\
RF & environment\_only & (3,14,7) & 171 & 0.007 & 0.224 & 0.000 & 0.988 & 0.000 & 0.000 & 0.530 & 0.015 \\
RF & no\_environment & (3,14,7) & 315 & 0.006 & 0.137 & 0.000 & 0.988 & 0.000 & 0.000 & 0.592 & 0.017 \\
RF & no\_peakflow & (3,14,7) & 355 & 0.023 & 0.765 & 0.000 & 0.988 & 0.000 & 0.000 & 0.493 & 0.014 \\
RF & no\_smartinhaler & (3,14,7) & 400 & 1.000 & 1.000 & 0.400 & 0.988 & 0.000 & 0.000 & 0.396 & 0.012 \\
RF & no\_smartwatch & (3,14,7) & 270 & 0.007 & 0.251 & 0.000 & 0.967 & 0.111 & 0.250 & 0.395 & 0.133 \\
RF & peakflow\_only & (3,14,7) & 131 & 0.006 & 0.126 & 0.000 & 0.722 & 0.011 & 0.250 & 0.274 & 0.030 \\
RF & smartinhaler\_only & (3,14,7) & 86 & 0.011 & 0.552 & 0.000 & 0.882 & 0.000 & 0.000 & 0.686 & 0.023 \\
RF & smartwatch\_only & (3,14,7) & 216 & 0.006 & 0.087 & 0.000 & 0.988 & 0.000 & 0.000 & 0.795 & 0.050 \\
RF & all & (3,14,14) & 427 & 0.021 & 0.749 & 0.000 & 0.957 & 0.000 & 0.000 & 0.809 & 0.039 \\
RF & environment\_only & (3,14,14) & 171 & 0.007 & 0.202 & 0.000 & 0.988 & 0.000 & 0.000 & 0.720 & 0.029 \\
RF & no\_environment & (3,14,14) & 315 & 0.005 & 0.000 & 0.000 & 0.960 & 0.154 & 0.500 & 0.916 & 0.356 \\
RF & no\_peakflow & (3,14,14) & 355 & 0.009 & 0.415 & 0.000 & 0.988 & 0.000 & 0.000 & 0.878 & 0.332 \\
RF & no\_smartinhaler & (3,14,14) & 400 & 0.011 & 0.525 & 0.000 & 0.988 & 0.000 & 0.000 & 0.855 & 0.058 \\
RF & no\_smartwatch & (3,14,14) & 270 & 0.007 & 0.240 & 0.000 & 0.901 & 0.033 & 0.250 & 0.441 & 0.024 \\
RF & peakflow\_only & (3,14,14) & 131 & 0.006 & 0.060 & 0.000 & 0.392 & 0.005 & 0.250 & 0.303 & 0.016 \\
RF & smartinhaler\_only & (3,14,14) & 86 & 0.006 & 0.131 & 0.000 & 0.975 & 0.167 & 0.250 & 0.745 & 0.062 \\
RF & smartwatch\_only & (3,14,14) & 216 & 0.005 & 0.000 & 0.000 & 0.966 & 0.111 & 0.250 & 0.889 & 0.166 \\
RF & all & (4,3,0) & 427 & 0.036 & 0.894 & 0.000 & 0.988 & 0.000 & 0.000 & 0.423 & 0.013 \\
RF & environment\_only & (4,3,0) & 171 & 0.008 & 0.472 & 0.000 & 0.988 & 0.000 & 0.000 & 0.558 & 0.024 \\
RF & no\_environment & (4,3,0) & 315 & 0.007 & 0.329 & 0.000 & 0.988 & 0.000 & 0.000 & 0.340 & 0.010 \\
RF & no\_peakflow & (4,3,0) & 355 & 0.008 & 0.472 & 0.000 & 0.988 & 0.000 & 0.000 & 0.561 & 0.018 \\
RF & no\_smartinhaler & (4,3,0) & 400 & 0.019 & 0.778 & 0.000 & 0.988 & 0.000 & 0.000 & 0.407 & 0.012 \\
RF & no\_smartwatch & (4,3,0) & 270 & 0.015 & 0.722 & 0.000 & 0.988 & 0.000 & 0.000 & 0.435 & 0.016 \\
RF & peakflow\_only & (4,3,0) & 131 & 0.017 & 0.822 & 0.000 & 0.988 & 0.000 & 0.000 & 0.498 & 0.014 \\
RF & smartinhaler\_only & (4,3,0) & 86 & 0.007 & 0.412 & 0.000 & 0.988 & 0.000 & 0.000 & 0.256 & 0.009 \\
RF & smartwatch\_only & (4,3,0) & 216 & 0.013 & 0.657 & 0.000 & 0.988 & 0.000 & 0.000 & 0.328 & 0.010 \\
RF & all & (4,3,7) & 427 & 0.016 & 0.727 & 0.000 & 0.987 & 0.000 & 0.000 & 0.521 & 0.019 \\
RF & environment\_only & (4,3,7) & 171 & 0.010 & 0.546 & 0.000 & 0.987 & 0.000 & 0.000 & 0.539 & 0.021 \\
RF & no\_environment & (4,3,7) & 315 & 0.007 & 0.329 & 0.000 & 0.987 & 0.000 & 0.000 & 0.275 & 0.010 \\
RF & no\_peakflow & (4,3,7) & 355 & 0.011 & 0.593 & 0.000 & 0.987 & 0.000 & 0.000 & 0.481 & 0.016 \\
RF & no\_smartinhaler & (4,3,7) & 400 & 0.014 & 0.690 & 0.000 & 0.987 & 0.000 & 0.000 & 0.531 & 0.018 \\
RF & no\_smartwatch & (4,3,7) & 270 & 0.009 & 0.514 & 0.000 & 0.987 & 0.000 & 0.000 & 0.452 & 0.018 \\
RF & peakflow\_only & (4,3,7) & 131 & 0.011 & 0.669 & 0.000 & 0.987 & 0.000 & 0.000 & 0.348 & 0.011 \\
RF & smartinhaler\_only & (4,3,7) & 86 & 0.006 & 0.366 & 0.000 & 0.987 & 0.000 & 0.000 & 0.250 & 0.010 \\
RF & smartwatch\_only & (4,3,7) & 216 & 0.010 & 0.583 & 0.000 & 0.987 & 0.000 & 0.000 & 0.425 & 0.012 \\
RF & all & (4,3,14) & 427 & 0.011 & 0.583 & 0.000 & 0.987 & 0.000 & 0.000 & 0.467 & 0.014 \\
RF & environment\_only & (4,3,14) & 171 & 0.006 & 0.245 & 0.000 & 0.987 & 0.000 & 0.000 & 0.569 & 0.026 \\
RF & no\_environment & (4,3,14) & 315 & 0.006 & 0.199 & 0.000 & 0.987 & 0.000 & 0.000 & 0.314 & 0.011 \\
RF & no\_peakflow & (4,3,14) & 355 & 0.007 & 0.370 & 0.000 & 0.987 & 0.000 & 0.000 & 0.587 & 0.035 \\
RF & no\_smartinhaler & (4,3,14) & 400 & 0.008 & 0.463 & 0.000 & 0.987 & 0.000 & 0.000 & 0.467 & 0.025 \\
RF & no\_smartwatch & (4,3,14) & 270 & 0.005 & 0.181 & 0.000 & 0.987 & 0.000 & 0.000 & 0.487 & 0.019 \\
RF & peakflow\_only & (4,3,14) & 131 & 0.006 & 0.350 & 0.000 & 0.987 & 0.000 & 0.000 & 0.388 & 0.012 \\
RF & smartinhaler\_only & (4,3,14) & 86 & 0.006 & 0.315 & 0.000 & 0.987 & 0.000 & 0.000 & 0.239 & 0.010 \\
RF & smartwatch\_only & (4,3,14) & 216 & 0.010 & 0.546 & 0.000 & 0.987 & 0.000 & 0.000 & 0.542 & 0.017 \\
RF & all & (4,7,0) & 427 & 0.016 & 0.708 & 0.000 & 0.987 & 0.000 & 0.000 & 0.636 & 0.235 \\
RF & environment\_only & (4,7,0) & 171 & 0.013 & 0.629 & 0.000 & 0.987 & 0.000 & 0.000 & 0.574 & 0.059 \\
RF & no\_environment & (4,7,0) & 315 & 0.009 & 0.470 & 0.000 & 0.984 & 0.000 & 0.000 & 0.540 & 0.017 \\
RF & no\_peakflow & (4,7,0) & 355 & 0.012 & 0.599 & 0.000 & 0.987 & 0.000 & 0.000 & 0.644 & 0.042 \\
RF & no\_smartinhaler & (4,7,0) & 400 & 0.010 & 0.535 & 0.000 & 0.987 & 0.000 & 0.000 & 0.578 & 0.048 \\
RF & no\_smartwatch & (4,7,0) & 270 & 0.014 & 0.649 & 0.000 & 0.987 & 0.000 & 0.000 & 0.531 & 0.023 \\
RF & peakflow\_only & (4,7,0) & 131 & 0.013 & 0.693 & 0.000 & 0.987 & 0.000 & 0.000 & 0.310 & 0.037 \\
RF & smartinhaler\_only & (4,7,0) & 86 & 0.012 & 0.636 & 0.000 & 0.987 & 0.000 & 0.000 & 0.691 & 0.028 \\
RF & smartwatch\_only & (4,7,0) & 216 & 0.009 & 0.450 & 0.000 & 0.982 & 0.000 & 0.000 & 0.566 & 0.064 \\
RF & all & (4,7,7) & 427 & 0.010 & 0.530 & 0.000 & 0.987 & 0.000 & 0.000 & 0.517 & 0.026 \\
RF & environment\_only & (4,7,7) & 171 & 0.006 & 0.168 & 0.000 & 0.987 & 0.000 & 0.000 & 0.577 & 0.055 \\
RF & no\_environment & (4,7,7) & 315 & 0.008 & 0.386 & 0.000 & 0.987 & 0.000 & 0.000 & 0.311 & 0.011 \\
RF & no\_peakflow & (4,7,7) & 355 & 0.010 & 0.525 & 0.000 & 0.987 & 0.000 & 0.000 & 0.701 & 0.103 \\
RF & no\_smartinhaler & (4,7,7) & 400 & 0.009 & 0.450 & 0.000 & 0.987 & 0.000 & 0.000 & 0.499 & 0.023 \\
RF & no\_smartwatch & (4,7,7) & 270 & 0.007 & 0.337 & 0.000 & 0.987 & 0.000 & 0.000 & 0.529 & 0.018 \\
RF & peakflow\_only & (4,7,7) & 131 & 0.019 & 0.802 & 0.000 & 0.987 & 0.000 & 0.000 & 0.178 & 0.010 \\
RF & smartinhaler\_only & (4,7,7) & 86 & 0.013 & 0.686 & 0.000 & 0.987 & 0.000 & 0.000 & 0.516 & 0.017 \\
RF & smartwatch\_only & (4,7,7) & 216 & 0.011 & 0.574 & 0.000 & 0.987 & 0.000 & 0.000 & 0.618 & 0.051 \\
RF & all & (4,7,14) & 427 & 0.009 & 0.460 & 0.000 & 0.986 & 0.000 & 0.000 & 0.529 & 0.084 \\
RF & environment\_only & (4,7,14) & 171 & 0.007 & 0.277 & 0.000 & 0.975 & 0.000 & 0.000 & 0.633 & 0.060 \\
RF & no\_environment & (4,7,14) & 315 & 0.007 & 0.351 & 0.000 & 0.986 & 0.000 & 0.000 & 0.467 & 0.015 \\
RF & no\_peakflow & (4,7,14) & 355 & 0.017 & 0.723 & 0.000 & 0.981 & 0.000 & 0.000 & 0.720 & 0.183 \\
RF & no\_smartinhaler & (4,7,14) & 400 & 0.010 & 0.500 & 0.000 & 0.986 & 0.000 & 0.000 & 0.565 & 0.092 \\
RF & no\_smartwatch & (4,7,14) & 270 & 0.005 & 0.069 & 0.000 & 0.986 & 0.000 & 0.000 & 0.434 & 0.015 \\
RF & peakflow\_only & (4,7,14) & 131 & 0.009 & 0.505 & 0.000 & 0.972 & 0.000 & 0.000 & 0.180 & 0.010 \\
RF & smartinhaler\_only & (4,7,14) & 86 & 0.013 & 0.661 & 0.000 & 0.986 & 0.000 & 0.000 & 0.737 & 0.031 \\
RF & smartwatch\_only & (4,7,14) & 216 & 0.009 & 0.475 & 0.000 & 0.986 & 0.000 & 0.000 & 0.719 & 0.130 \\
RF & all & (4,14,0) & 427 & 0.045 & 0.885 & 0.000 & 0.989 & 0.000 & 0.000 & 0.595 & 0.035 \\
RF & environment\_only & (4,14,0) & 171 & 0.011 & 0.514 & 0.000 & 0.989 & 0.000 & 0.000 & 0.504 & 0.025 \\
RF & no\_environment & (4,14,0) & 315 & 0.006 & 0.142 & 0.000 & 0.989 & 0.000 & 0.000 & 0.677 & 0.024 \\
RF & no\_peakflow & (4,14,0) & 355 & 0.021 & 0.749 & 0.000 & 0.989 & 0.000 & 0.000 & 0.754 & 0.059 \\
RF & no\_smartinhaler & (4,14,0) & 400 & 0.023 & 0.765 & 0.000 & 0.971 & 0.000 & 0.000 & 0.652 & 0.034 \\
RF & no\_smartwatch & (4,14,0) & 270 & 0.010 & 0.448 & 0.000 & 0.986 & 0.000 & 0.000 & 0.520 & 0.015 \\
RF & peakflow\_only & (4,14,0) & 131 & 0.008 & 0.311 & 0.000 & 0.843 & 0.019 & 0.250 & 0.291 & 0.038 \\
RF & smartinhaler\_only & (4,14,0) & 86 & 0.006 & 0.142 & 0.000 & 0.989 & 0.000 & 0.000 & 0.753 & 0.027 \\
RF & smartwatch\_only & (4,14,0) & 216 & 0.006 & 0.087 & 0.000 & 0.989 & 0.000 & 0.000 & 0.727 & 0.054 \\
RF & all & (4,14,7) & 427 & 0.250 & 0.989 & 0.000 & 0.988 & 0.000 & 0.000 & 0.664 & 0.031 \\
RF & environment\_only & (4,14,7) & 171 & 0.007 & 0.186 & 0.000 & 0.988 & 0.000 & 0.000 & 0.443 & 0.014 \\
RF & no\_environment & (4,14,7) & 315 & 0.006 & 0.115 & 0.000 & 0.988 & 0.000 & 0.000 & 0.653 & 0.029 \\
RF & no\_peakflow & (4,14,7) & 355 & 0.018 & 0.699 & 0.000 & 0.988 & 0.000 & 0.000 & 0.694 & 0.052 \\
RF & no\_smartinhaler & (4,14,7) & 400 & 0.067 & 0.926 & 0.000 & 0.988 & 0.000 & 0.000 & 0.654 & 0.030 \\
RF & no\_smartwatch & (4,14,7) & 270 & 0.014 & 0.642 & 0.000 & 0.988 & 0.000 & 0.000 & 0.504 & 0.014 \\
RF & peakflow\_only & (4,14,7) & 131 & 0.011 & 0.541 & 0.000 & 0.873 & 0.000 & 0.000 & 0.336 & 0.014 \\
RF & smartinhaler\_only & (4,14,7) & 86 & 0.011 & 0.557 & 0.000 & 0.988 & 0.000 & 0.000 & 0.522 & 0.017 \\
RF & smartwatch\_only & (4,14,7) & 216 & 0.006 & 0.123 & 0.000 & 0.988 & 0.000 & 0.000 & 0.764 & 0.048 \\
RF & all & (4,14,14) & 427 & 0.015 & 0.650 & 0.000 & 0.991 & 1.000 & 0.250 & 0.819 & 0.295 \\
RF & environment\_only & (4,14,14) & 171 & 0.008 & 0.295 & 0.000 & 0.988 & 0.000 & 0.000 & 0.597 & 0.263 \\
RF & no\_environment & (4,14,14) & 315 & 0.006 & 0.104 & 0.000 & 0.975 & 0.000 & 0.000 & 0.726 & 0.054 \\
RF & no\_peakflow & (4,14,14) & 355 & 0.014 & 0.612 & 0.000 & 0.991 & 1.000 & 0.250 & 0.830 & 0.371 \\
RF & no\_smartinhaler & (4,14,14) & 400 & 0.020 & 0.727 & 0.000 & 0.978 & 0.200 & 0.250 & 0.695 & 0.076 \\
RF & no\_smartwatch & (4,14,14) & 270 & 0.008 & 0.295 & 0.000 & 0.892 & 0.000 & 0.000 & 0.493 & 0.016 \\
RF & peakflow\_only & (4,14,14) & 131 & 0.006 & 0.082 & 0.000 & 0.751 & 0.013 & 0.250 & 0.322 & 0.017 \\
RF & smartinhaler\_only & (4,14,14) & 86 & 0.006 & 0.142 & 0.000 & 0.972 & 0.143 & 0.250 & 0.735 & 0.069 \\
RF & smartwatch\_only & (4,14,14) & 216 & 0.006 & 0.120 & 0.000 & 0.991 & 1.000 & 0.250 & 0.734 & 0.531 \\
XGB & all & (2,3,0) & 427 & 0.007 & 0.347 & 0.000 & 0.990 & 0.000 & 0.000 & 0.464 & 0.009 \\
XGB & environment\_only & (2,3,0) & 171 & 0.006 & 0.306 & 0.000 & 0.993 & 0.000 & 0.000 & 0.568 & 0.016 \\
XGB & no\_environment & (2,3,0) & 315 & 0.006 & 0.301 & 0.000 & 0.949 & 0.000 & 0.000 & 0.404 & 0.008 \\
XGB & no\_peakflow & (2,3,0) & 355 & 0.007 & 0.347 & 0.000 & 0.993 & 0.000 & 0.000 & 0.670 & 0.027 \\
XGB & no\_smartinhaler & (2,3,0) & 400 & 0.007 & 0.347 & 0.000 & 0.993 & 0.000 & 0.000 & 0.490 & 0.009 \\
XGB & no\_smartwatch & (2,3,0) & 270 & 0.009 & 0.528 & 0.000 & 0.980 & 0.000 & 0.000 & 0.485 & 0.010 \\
XGB & peakflow\_only & (2,3,0) & 131 & 0.009 & 0.581 & 0.000 & 0.919 & 0.000 & 0.000 & 0.354 & 0.007 \\
XGB & smartinhaler\_only & (2,3,0) & 86 & 0.024 & 0.889 & 0.000 & 0.993 & 0.000 & 0.000 & 0.705 & 0.092 \\
XGB & smartwatch\_only & (2,3,0) & 216 & 0.007 & 0.375 & 0.000 & 0.993 & 0.000 & 0.000 & 0.652 & 0.026 \\
XGB & all & (2,3,7) & 427 & 0.006 & 0.185 & 0.000 & 0.980 & 0.000 & 0.000 & 0.510 & 0.010 \\
XGB & environment\_only & (2,3,7) & 171 & 0.007 & 0.315 & 0.000 & 0.992 & 0.000 & 0.000 & 0.523 & 0.013 \\
XGB & no\_environment & (2,3,7) & 315 & 0.005 & 0.167 & 0.000 & 0.721 & 0.000 & 0.000 & 0.443 & 0.009 \\
XGB & no\_peakflow & (2,3,7) & 355 & 0.005 & 0.162 & 0.000 & 0.992 & 0.000 & 0.000 & 0.659 & 0.017 \\
XGB & no\_smartinhaler & (2,3,7) & 400 & 0.005 & 0.162 & 0.000 & 0.975 & 0.000 & 0.000 & 0.524 & 0.010 \\
XGB & no\_smartwatch & (2,3,7) & 270 & 0.007 & 0.347 & 0.000 & 0.952 & 0.000 & 0.000 & 0.446 & 0.010 \\
XGB & peakflow\_only & (2,3,7) & 131 & 0.008 & 0.539 & 0.000 & 0.744 & 0.000 & 0.000 & 0.286 & 0.007 \\
XGB & smartinhaler\_only & (2,3,7) & 86 & 0.020 & 0.856 & 0.000 & 0.992 & 0.000 & 0.000 & 0.665 & 0.031 \\
XGB & smartwatch\_only & (2,3,7) & 216 & 0.005 & 0.157 & 0.000 & 0.992 & 0.000 & 0.000 & 0.610 & 0.043 \\
XGB & all & (2,3,14) & 427 & 0.005 & 0.116 & 0.000 & 0.945 & 0.000 & 0.000 & 0.493 & 0.010 \\
XGB & environment\_only & (2,3,14) & 171 & 0.006 & 0.310 & 0.000 & 0.992 & 0.000 & 0.000 & 0.632 & 0.019 \\
XGB & no\_environment & (2,3,14) & 315 & 0.005 & 0.111 & 0.000 & 0.830 & 0.000 & 0.000 & 0.461 & 0.010 \\
XGB & no\_peakflow & (2,3,14) & 355 & 0.005 & 0.088 & 0.000 & 0.992 & 0.000 & 0.000 & 0.751 & 0.025 \\
XGB & no\_smartinhaler & (2,3,14) & 400 & 0.005 & 0.111 & 0.000 & 0.958 & 0.000 & 0.000 & 0.455 & 0.009 \\
XGB & no\_smartwatch & (2,3,14) & 270 & 0.008 & 0.454 & 0.000 & 0.882 & 0.000 & 0.000 & 0.404 & 0.009 \\
XGB & peakflow\_only & (2,3,14) & 131 & 0.008 & 0.475 & 0.000 & 0.686 & 0.000 & 0.000 & 0.260 & 0.007 \\
XGB & smartinhaler\_only & (2,3,14) & 86 & 0.025 & 0.898 & 0.000 & 0.992 & 0.000 & 0.000 & 0.652 & 0.025 \\
XGB & smartwatch\_only & (2,3,14) & 216 & 0.005 & 0.130 & 0.000 & 0.992 & 0.000 & 0.000 & 0.664 & 0.061 \\
XGB & all & (2,7,0) & 427 & 0.006 & 0.178 & 0.000 & 0.990 & 0.000 & 0.000 & 0.825 & 0.102 \\
XGB & environment\_only & (2,7,0) & 171 & 0.008 & 0.416 & 0.000 & 0.966 & 0.000 & 0.000 & 0.584 & 0.038 \\
XGB & no\_environment & (2,7,0) & 315 & 0.006 & 0.238 & 0.000 & 0.979 & 0.000 & 0.000 & 0.540 & 0.011 \\
XGB & no\_peakflow & (2,7,0) & 355 & 0.006 & 0.233 & 0.000 & 0.995 & 1.000 & 0.333 & 0.784 & 0.371 \\
XGB & no\_smartinhaler & (2,7,0) & 400 & 0.006 & 0.218 & 0.000 & 0.990 & 0.000 & 0.000 & 0.770 & 0.180 \\
XGB & no\_smartwatch & (2,7,0) & 270 & 0.010 & 0.530 & 0.000 & 0.938 & 0.000 & 0.000 & 0.494 & 0.011 \\
XGB & peakflow\_only & (2,7,0) & 131 & 0.016 & 0.752 & 0.000 & 0.599 & 0.000 & 0.000 & 0.278 & 0.007 \\
XGB & smartinhaler\_only & (2,7,0) & 86 & 0.036 & 0.918 & 0.000 & 0.927 & 0.037 & 0.333 & 0.539 & 0.090 \\
XGB & smartwatch\_only & (2,7,0) & 216 & 0.007 & 0.332 & 0.000 & 0.984 & 0.000 & 0.000 & 0.882 & 0.095 \\
XGB & all & (2,7,7) & 427 & 0.005 & 0.084 & 0.000 & 0.973 & 0.111 & 0.333 & 0.683 & 0.124 \\
XGB & environment\_only & (2,7,7) & 171 & 0.006 & 0.149 & 0.000 & 0.916 & 0.033 & 0.333 & 0.658 & 0.021 \\
XGB & no\_environment & (2,7,7) & 315 & 0.005 & 0.084 & 0.000 & 0.816 & 0.000 & 0.000 & 0.470 & 0.010 \\
XGB & no\_peakflow & (2,7,7) & 355 & 0.005 & 0.045 & 0.000 & 0.995 & 1.000 & 0.333 & 0.856 & 0.354 \\
XGB & no\_smartinhaler & (2,7,7) & 400 & 0.005 & 0.084 & 0.000 & 0.981 & 0.167 & 0.333 & 0.679 & 0.177 \\
XGB & no\_smartwatch & (2,7,7) & 270 & 0.006 & 0.124 & 0.000 & 0.792 & 0.000 & 0.000 & 0.571 & 0.015 \\
XGB & peakflow\_only & (2,7,7) & 131 & 0.010 & 0.574 & 0.000 & 0.470 & 0.005 & 0.333 & 0.218 & 0.007 \\
XGB & smartinhaler\_only & (2,7,7) & 86 & 0.040 & 0.933 & 0.000 & 0.957 & 0.067 & 0.333 & 0.757 & 0.188 \\
XGB & smartwatch\_only & (2,7,7) & 216 & 0.007 & 0.282 & 0.000 & 0.970 & 0.100 & 0.333 & 0.727 & 0.397 \\
XGB & all & (2,7,14) & 427 & 0.005 & 0.020 & 0.000 & 0.983 & 0.200 & 0.333 & 0.613 & 0.342 \\
XGB & environment\_only & (2,7,14) & 171 & 0.006 & 0.178 & 0.000 & 0.966 & 0.091 & 0.333 & 0.723 & 0.069 \\
XGB & no\_environment & (2,7,14) & 315 & 0.005 & 0.089 & 0.000 & 0.860 & 0.000 & 0.000 & 0.515 & 0.012 \\
XGB & no\_peakflow & (2,7,14) & 355 & 0.005 & 0.030 & 0.000 & 0.994 & 1.000 & 0.333 & 0.735 & 0.345 \\
XGB & no\_smartinhaler & (2,7,14) & 400 & 0.005 & 0.020 & 0.000 & 0.992 & 0.500 & 0.333 & 0.675 & 0.343 \\
XGB & no\_smartwatch & (2,7,14) & 270 & 0.006 & 0.114 & 0.000 & 0.852 & 0.019 & 0.333 & 0.508 & 0.073 \\
XGB & peakflow\_only & (2,7,14) & 131 & 0.009 & 0.525 & 0.000 & 0.617 & 0.000 & 0.000 & 0.234 & 0.007 \\
XGB & smartinhaler\_only & (2,7,14) & 86 & 0.031 & 0.899 & 0.000 & 0.975 & 0.125 & 0.333 & 0.767 & 0.085 \\
XGB & smartwatch\_only & (2,7,14) & 216 & 0.006 & 0.183 & 0.000 & 0.992 & 0.000 & 0.000 & 0.861 & 0.201 \\
XGB & all & (2,14,0) & 427 & 0.009 & 0.421 & 0.000 & 0.983 & 0.000 & 0.000 & 0.698 & 0.018 \\
XGB & environment\_only & (2,14,0) & 171 & 0.006 & 0.148 & 0.000 & 0.954 & 0.000 & 0.000 & 0.556 & 0.009 \\
XGB & no\_environment & (2,14,0) & 315 & 0.005 & 0.000 & 0.000 & 0.977 & 0.000 & 0.000 & 0.777 & 0.029 \\
XGB & no\_peakflow & (2,14,0) & 355 & 0.009 & 0.372 & 0.000 & 0.986 & 0.000 & 0.000 & 0.653 & 0.014 \\
XGB & no\_smartinhaler & (2,14,0) & 400 & 0.009 & 0.393 & 0.000 & 0.968 & 0.000 & 0.000 & 0.640 & 0.013 \\
XGB & no\_smartwatch & (2,14,0) & 270 & 0.007 & 0.224 & 0.000 & 0.922 & 0.000 & 0.000 & 0.630 & 0.024 \\
XGB & peakflow\_only & (2,14,0) & 131 & 0.006 & 0.137 & 0.000 & 0.374 & 0.005 & 0.500 & 0.503 & 0.034 \\
XGB & smartinhaler\_only & (2,14,0) & 86 & 0.045 & 0.923 & 0.000 & 0.856 & 0.020 & 0.500 & 0.741 & 0.024 \\
XGB & smartwatch\_only & (2,14,0) & 216 & 0.005 & 0.000 & 0.000 & 0.960 & 0.000 & 0.000 & 0.660 & 0.017 \\
XGB & all & (2,14,7) & 427 & 0.007 & 0.262 & 0.000 & 0.988 & 0.000 & 0.000 & 0.620 & 0.013 \\
XGB & environment\_only & (2,14,7) & 171 & 0.007 & 0.208 & 0.000 & 0.929 & 0.000 & 0.000 & 0.461 & 0.008 \\
XGB & no\_environment & (2,14,7) & 315 & 0.005 & 0.000 & 0.000 & 0.926 & 0.000 & 0.000 & 0.627 & 0.013 \\
XGB & no\_peakflow & (2,14,7) & 355 & 0.007 & 0.257 & 0.000 & 0.985 & 0.000 & 0.000 & 0.811 & 0.023 \\
XGB & no\_smartinhaler & (2,14,7) & 400 & 0.007 & 0.197 & 0.000 & 0.961 & 0.000 & 0.000 & 0.536 & 0.010 \\
XGB & no\_smartwatch & (2,14,7) & 270 & 0.007 & 0.235 & 0.000 & 0.732 & 0.000 & 0.000 & 0.407 & 0.007 \\
XGB & peakflow\_only & (2,14,7) & 131 & 0.006 & 0.164 & 0.000 & 0.235 & 0.004 & 0.500 & 0.425 & 0.013 \\
XGB & smartinhaler\_only & (2,14,7) & 86 & 0.056 & 0.945 & 0.000 & 0.824 & 0.000 & 0.000 & 0.696 & 0.014 \\
XGB & smartwatch\_only & (2,14,7) & 216 & 0.006 & 0.022 & 0.000 & 0.985 & 0.000 & 0.000 & 0.696 & 0.015 \\
XGB & all & (2,14,14) & 427 & 0.007 & 0.257 & 0.000 & 0.904 & 0.032 & 0.500 & 0.819 & 0.041 \\
XGB & environment\_only & (2,14,14) & 171 & 0.007 & 0.273 & 0.000 & 0.882 & 0.000 & 0.000 & 0.547 & 0.017 \\
XGB & no\_environment & (2,14,14) & 315 & 0.005 & 0.000 & 0.000 & 0.969 & 0.100 & 0.500 & 0.880 & 0.085 \\
XGB & no\_peakflow & (2,14,14) & 355 & 0.006 & 0.115 & 0.000 & 0.960 & 0.077 & 0.500 & 0.808 & 0.054 \\
XGB & no\_smartinhaler & (2,14,14) & 400 & 0.006 & 0.104 & 0.000 & 0.910 & 0.034 & 0.500 & 0.812 & 0.034 \\
XGB & no\_smartwatch & (2,14,14) & 270 & 0.007 & 0.273 & 0.000 & 0.624 & 0.000 & 0.000 & 0.441 & 0.008 \\
XGB & peakflow\_only & (2,14,14) & 131 & 0.005 & 0.033 & 0.000 & 0.357 & 0.005 & 0.500 & 0.305 & 0.007 \\
XGB & smartinhaler\_only & (2,14,14) & 86 & 0.007 & 0.257 & 0.000 & 0.894 & 0.056 & 1.000 & 0.912 & 0.050 \\
XGB & smartwatch\_only & (2,14,14) & 216 & 0.005 & 0.005 & 0.000 & 0.941 & 0.053 & 0.500 & 0.872 & 0.137 \\
XGB & all & (3,3,0) & 427 & 0.006 & 0.264 & 0.000 & 0.983 & 0.000 & 0.000 & 0.279 & 0.011 \\
XGB & environment\_only & (3,3,0) & 171 & 0.008 & 0.412 & 0.000 & 0.988 & 0.000 & 0.000 & 0.399 & 0.012 \\
XGB & no\_environment & (3,3,0) & 315 & 0.006 & 0.231 & 0.000 & 0.978 & 0.000 & 0.000 & 0.305 & 0.010 \\
XGB & no\_peakflow & (3,3,0) & 355 & 0.006 & 0.278 & 0.000 & 0.988 & 0.000 & 0.000 & 0.377 & 0.011 \\
XGB & no\_smartinhaler & (3,3,0) & 400 & 0.006 & 0.301 & 0.000 & 0.983 & 0.000 & 0.000 & 0.283 & 0.011 \\
XGB & no\_smartwatch & (3,3,0) & 270 & 0.009 & 0.491 & 0.000 & 0.968 & 0.000 & 0.000 & 0.210 & 0.009 \\
XGB & peakflow\_only & (3,3,0) & 131 & 0.010 & 0.637 & 0.000 & 0.854 & 0.000 & 0.000 & 0.322 & 0.010 \\
XGB & smartinhaler\_only & (3,3,0) & 86 & 0.006 & 0.287 & 0.000 & 0.988 & 0.000 & 0.000 & 0.376 & 0.011 \\
XGB & smartwatch\_only & (3,3,0) & 216 & 0.006 & 0.236 & 0.000 & 0.988 & 0.000 & 0.000 & 0.456 & 0.018 \\
XGB & all & (3,3,7) & 427 & 0.005 & 0.125 & 0.000 & 0.975 & 0.000 & 0.000 & 0.345 & 0.011 \\
XGB & environment\_only & (3,3,7) & 171 & 0.008 & 0.468 & 0.000 & 0.987 & 0.000 & 0.000 & 0.528 & 0.020 \\
XGB & no\_environment & (3,3,7) & 315 & 0.005 & 0.120 & 0.000 & 0.932 & 0.000 & 0.000 & 0.332 & 0.011 \\
XGB & no\_peakflow & (3,3,7) & 355 & 0.005 & 0.116 & 0.000 & 0.987 & 0.000 & 0.000 & 0.576 & 0.019 \\
XGB & no\_smartinhaler & (3,3,7) & 400 & 0.005 & 0.162 & 0.000 & 0.970 & 0.000 & 0.000 & 0.320 & 0.011 \\
XGB & no\_smartwatch & (3,3,7) & 270 & 0.008 & 0.417 & 0.000 & 0.955 & 0.000 & 0.000 & 0.389 & 0.012 \\
XGB & peakflow\_only & (3,3,7) & 131 & 0.008 & 0.525 & 0.000 & 0.876 & 0.000 & 0.000 & 0.302 & 0.010 \\
XGB & smartinhaler\_only & (3,3,7) & 86 & 0.027 & 0.912 & 0.000 & 0.987 & 0.000 & 0.000 & 0.384 & 0.012 \\
XGB & smartwatch\_only & (3,3,7) & 216 & 0.005 & 0.139 & 0.000 & 0.987 & 0.000 & 0.000 & 0.361 & 0.019 \\
XGB & all & (3,3,14) & 427 & 0.005 & 0.111 & 0.000 & 0.977 & 0.000 & 0.000 & 0.461 & 0.016 \\
XGB & environment\_only & (3,3,14) & 171 & 0.005 & 0.181 & 0.000 & 0.987 & 0.000 & 0.000 & 0.470 & 0.022 \\
XGB & no\_environment & (3,3,14) & 315 & 0.005 & 0.037 & 0.000 & 0.862 & 0.000 & 0.000 & 0.365 & 0.012 \\
XGB & no\_peakflow & (3,3,14) & 355 & 0.005 & 0.060 & 0.000 & 0.987 & 0.000 & 0.000 & 0.672 & 0.043 \\
XGB & no\_smartinhaler & (3,3,14) & 400 & 0.005 & 0.042 & 0.000 & 0.974 & 0.000 & 0.000 & 0.522 & 0.016 \\
XGB & no\_smartwatch & (3,3,14) & 270 & 0.005 & 0.162 & 0.000 & 0.961 & 0.000 & 0.000 & 0.323 & 0.011 \\
XGB & peakflow\_only & (3,3,14) & 131 & 0.007 & 0.382 & 0.000 & 0.776 & 0.000 & 0.000 & 0.362 & 0.012 \\
XGB & smartinhaler\_only & (3,3,14) & 86 & 0.005 & 0.093 & 0.000 & 0.987 & 0.000 & 0.000 & 0.336 & 0.011 \\
XGB & smartwatch\_only & (3,3,14) & 216 & 0.005 & 0.069 & 0.000 & 0.987 & 0.000 & 0.000 & 0.704 & 0.045 \\
XGB & all & (3,7,0) & 427 & 0.006 & 0.193 & 0.000 & 0.982 & 0.250 & 0.200 & 0.677 & 0.218 \\
XGB & environment\_only & (3,7,0) & 171 & 0.006 & 0.178 & 0.000 & 0.959 & 0.000 & 0.000 & 0.289 & 0.010 \\
XGB & no\_environment & (3,7,0) & 315 & 0.006 & 0.228 & 0.000 & 0.904 & 0.000 & 0.000 & 0.404 & 0.012 \\
XGB & no\_peakflow & (3,7,0) & 355 & 0.007 & 0.282 & 0.000 & 0.984 & 0.333 & 0.200 & 0.606 & 0.116 \\
XGB & no\_smartinhaler & (3,7,0) & 400 & 0.006 & 0.153 & 0.000 & 0.987 & 0.500 & 0.200 & 0.651 & 0.217 \\
XGB & no\_smartwatch & (3,7,0) & 270 & 0.006 & 0.173 & 0.000 & 0.935 & 0.000 & 0.000 & 0.310 & 0.011 \\
XGB & peakflow\_only & (3,7,0) & 131 & 0.015 & 0.733 & 0.000 & 0.573 & 0.006 & 0.200 & 0.304 & 0.011 \\
XGB & smartinhaler\_only & (3,7,0) & 86 & 0.030 & 0.894 & 0.000 & 0.886 & 0.047 & 0.400 & 0.632 & 0.044 \\
XGB & smartwatch\_only & (3,7,0) & 216 & 0.007 & 0.337 & 0.000 & 0.982 & 0.250 & 0.200 & 0.618 & 0.244 \\
XGB & all & (3,7,7) & 427 & 0.006 & 0.158 & 0.000 & 0.973 & 0.000 & 0.000 & 0.347 & 0.014 \\
XGB & environment\_only & (3,7,7) & 171 & 0.006 & 0.163 & 0.000 & 0.960 & 0.083 & 0.200 & 0.547 & 0.032 \\
XGB & no\_environment & (3,7,7) & 315 & 0.005 & 0.074 & 0.000 & 0.769 & 0.000 & 0.000 & 0.338 & 0.012 \\
XGB & no\_peakflow & (3,7,7) & 355 & 0.005 & 0.059 & 0.000 & 0.984 & 0.000 & 0.000 & 0.537 & 0.113 \\
XGB & no\_smartinhaler & (3,7,7) & 400 & 0.006 & 0.139 & 0.000 & 0.965 & 0.000 & 0.000 & 0.429 & 0.021 \\
XGB & no\_smartwatch & (3,7,7) & 270 & 0.005 & 0.104 & 0.000 & 0.944 & 0.000 & 0.000 & 0.456 & 0.016 \\
XGB & peakflow\_only & (3,7,7) & 131 & 0.006 & 0.297 & 0.000 & 0.573 & 0.000 & 0.000 & 0.112 & 0.009 \\
XGB & smartinhaler\_only & (3,7,7) & 86 & 0.006 & 0.245 & 0.000 & 0.973 & 0.000 & 0.000 & 0.556 & 0.026 \\
XGB & smartwatch\_only & (3,7,7) & 216 & 0.006 & 0.252 & 0.000 & 0.989 & 1.000 & 0.200 & 0.623 & 0.252 \\
XGB & all & (3,7,14) & 427 & 0.005 & 0.054 & 0.000 & 0.947 & 0.062 & 0.200 & 0.589 & 0.115 \\
XGB & environment\_only & (3,7,14) & 171 & 0.006 & 0.119 & 0.000 & 0.981 & 0.250 & 0.200 & 0.549 & 0.214 \\
XGB & no\_environment & (3,7,14) & 315 & 0.005 & 0.074 & 0.000 & 0.753 & 0.000 & 0.000 & 0.531 & 0.017 \\
XGB & no\_peakflow & (3,7,14) & 355 & 0.005 & 0.035 & 0.000 & 0.986 & 0.500 & 0.200 & 0.793 & 0.249 \\
XGB & no\_smartinhaler & (3,7,14) & 400 & 0.005 & 0.020 & 0.000 & 0.953 & 0.071 & 0.200 & 0.669 & 0.219 \\
XGB & no\_smartwatch & (3,7,14) & 270 & 0.006 & 0.183 & 0.000 & 0.872 & 0.023 & 0.200 & 0.488 & 0.037 \\
XGB & peakflow\_only & (3,7,14) & 131 & 0.005 & 0.064 & 0.000 & 0.706 & 0.000 & 0.000 & 0.143 & 0.009 \\
XGB & smartinhaler\_only & (3,7,14) & 86 & 0.005 & 0.101 & 0.000 & 0.983 & 0.000 & 0.000 & 0.856 & 0.089 \\
XGB & smartwatch\_only & (3,7,14) & 216 & 0.005 & 0.099 & 0.000 & 0.983 & 0.333 & 0.200 & 0.741 & 0.239 \\
XGB & all & (3,14,0) & 427 & 0.029 & 0.820 & 0.026 & 0.900 & 0.000 & 0.000 & 0.486 & 0.013 \\
XGB & environment\_only & (3,14,0) & 171 & 0.008 & 0.344 & 0.000 & 0.949 & 0.000 & 0.000 & 0.454 & 0.013 \\
XGB & no\_environment & (3,14,0) & 315 & 0.005 & 0.000 & 0.000 & 0.966 & 0.000 & 0.000 & 0.691 & 0.023 \\
XGB & no\_peakflow & (3,14,0) & 355 & 0.010 & 0.443 & 0.000 & 0.963 & 0.000 & 0.000 & 0.449 & 0.012 \\
XGB & no\_smartinhaler & (3,14,0) & 400 & 0.009 & 0.426 & 0.000 & 0.943 & 0.000 & 0.000 & 0.441 & 0.012 \\
XGB & no\_smartwatch & (3,14,0) & 270 & 0.007 & 0.219 & 0.000 & 0.883 & 0.000 & 0.000 & 0.405 & 0.014 \\
XGB & peakflow\_only & (3,14,0) & 131 & 0.006 & 0.137 & 0.000 & 0.294 & 0.004 & 0.250 & 0.285 & 0.069 \\
XGB & smartinhaler\_only & (3,14,0) & 86 & 0.008 & 0.350 & 0.000 & 0.866 & 0.061 & 0.750 & 0.858 & 0.075 \\
XGB & smartwatch\_only & (3,14,0) & 216 & 0.005 & 0.000 & 0.000 & 0.954 & 0.000 & 0.000 & 0.837 & 0.046 \\
XGB & all & (3,14,7) & 427 & 0.021 & 0.749 & 0.021 & 0.876 & 0.000 & 0.000 & 0.557 & 0.016 \\
XGB & environment\_only & (3,14,7) & 171 & 0.010 & 0.443 & 0.000 & 0.846 & 0.000 & 0.000 & 0.338 & 0.011 \\
XGB & no\_environment & (3,14,7) & 315 & 0.005 & 0.000 & 0.000 & 0.879 & 0.000 & 0.000 & 0.607 & 0.018 \\
XGB & no\_peakflow & (3,14,7) & 355 & 0.010 & 0.459 & 0.000 & 0.935 & 0.000 & 0.000 & 0.676 & 0.028 \\
XGB & no\_smartinhaler & (3,14,7) & 400 & 0.010 & 0.437 & 0.000 & 0.876 & 0.000 & 0.000 & 0.511 & 0.014 \\
XGB & no\_smartwatch & (3,14,7) & 270 & 0.008 & 0.284 & 0.000 & 0.811 & 0.000 & 0.000 & 0.382 & 0.011 \\
XGB & peakflow\_only & (3,14,7) & 131 & 0.005 & 0.022 & 0.000 & 0.435 & 0.005 & 0.250 & 0.247 & 0.012 \\
XGB & smartinhaler\_only & (3,14,7) & 86 & 0.048 & 0.929 & 0.000 & 0.873 & 0.024 & 0.250 & 0.532 & 0.030 \\
XGB & smartwatch\_only & (3,14,7) & 216 & 0.005 & 0.005 & 0.000 & 0.959 & 0.143 & 0.500 & 0.805 & 0.093 \\
XGB & all & (3,14,14) & 427 & 0.010 & 0.459 & 0.000 & 0.917 & 0.000 & 0.000 & 0.785 & 0.035 \\
XGB & environment\_only & (3,14,14) & 171 & 0.009 & 0.410 & 0.000 & 0.944 & 0.000 & 0.000 & 0.677 & 0.030 \\
XGB & no\_environment & (3,14,14) & 315 & 0.005 & 0.011 & 0.000 & 0.985 & 0.000 & 0.000 & 0.916 & 0.112 \\
XGB & no\_peakflow & (3,14,14) & 355 & 0.007 & 0.219 & 0.000 & 0.994 & 1.000 & 0.500 & 0.900 & 0.571 \\
XGB & no\_smartinhaler & (3,14,14) & 400 & 0.014 & 0.601 & 0.000 & 0.910 & 0.000 & 0.000 & 0.739 & 0.032 \\
XGB & no\_smartwatch & (3,14,14) & 270 & 0.009 & 0.415 & 0.000 & 0.784 & 0.000 & 0.000 & 0.287 & 0.010 \\
XGB & peakflow\_only & (3,14,14) & 131 & 0.014 & 0.634 & 0.000 & 0.407 & 0.005 & 0.250 & 0.267 & 0.015 \\
XGB & smartinhaler\_only & (3,14,14) & 86 & 0.045 & 0.923 & 0.000 & 0.966 & 0.182 & 0.500 & 0.980 & 0.330 \\
XGB & smartwatch\_only & (3,14,14) & 216 & 0.005 & 0.000 & 0.000 & 0.994 & 1.000 & 0.500 & 0.880 & 0.631 \\
XGB & all & (4,3,0) & 427 & 0.007 & 0.333 & 0.000 & 0.944 & 0.000 & 0.000 & 0.262 & 0.009 \\
XGB & environment\_only & (4,3,0) & 171 & 0.008 & 0.458 & 0.000 & 0.973 & 0.000 & 0.000 & 0.443 & 0.013 \\
XGB & no\_environment & (4,3,0) & 315 & 0.007 & 0.361 & 0.000 & 0.929 & 0.000 & 0.000 & 0.287 & 0.010 \\
XGB & no\_peakflow & (4,3,0) & 355 & 0.006 & 0.301 & 0.000 & 0.980 & 0.000 & 0.000 & 0.346 & 0.011 \\
XGB & no\_smartinhaler & (4,3,0) & 400 & 0.007 & 0.343 & 0.000 & 0.968 & 0.000 & 0.000 & 0.320 & 0.010 \\
XGB & no\_smartwatch & (4,3,0) & 270 & 0.011 & 0.588 & 0.000 & 0.946 & 0.000 & 0.000 & 0.426 & 0.013 \\
XGB & peakflow\_only & (4,3,0) & 131 & 0.007 & 0.470 & 0.000 & 0.856 & 0.000 & 0.000 & 0.280 & 0.010 \\
XGB & smartinhaler\_only & (4,3,0) & 86 & 0.006 & 0.264 & 0.000 & 0.907 & 0.000 & 0.000 & 0.192 & 0.009 \\
XGB & smartwatch\_only & (4,3,0) & 216 & 0.006 & 0.306 & 0.000 & 0.976 & 0.000 & 0.000 & 0.209 & 0.009 \\
XGB & all & (4,3,7) & 427 & 0.006 & 0.199 & 0.000 & 0.977 & 0.000 & 0.000 & 0.322 & 0.010 \\
XGB & environment\_only & (4,3,7) & 171 & 0.006 & 0.231 & 0.000 & 0.987 & 0.000 & 0.000 & 0.296 & 0.010 \\
XGB & no\_environment & (4,3,7) & 315 & 0.005 & 0.176 & 0.000 & 0.924 & 0.000 & 0.000 & 0.352 & 0.011 \\
XGB & no\_peakflow & (4,3,7) & 355 & 0.005 & 0.148 & 0.000 & 0.987 & 0.000 & 0.000 & 0.439 & 0.013 \\
XGB & no\_smartinhaler & (4,3,7) & 400 & 0.006 & 0.282 & 0.000 & 0.982 & 0.000 & 0.000 & 0.395 & 0.012 \\
XGB & no\_smartwatch & (4,3,7) & 270 & 0.007 & 0.324 & 0.000 & 0.977 & 0.000 & 0.000 & 0.376 & 0.012 \\
XGB & peakflow\_only & (4,3,7) & 131 & 0.006 & 0.377 & 0.000 & 0.892 & 0.000 & 0.000 & 0.417 & 0.012 \\
XGB & smartinhaler\_only & (4,3,7) & 86 & 0.006 & 0.282 & 0.000 & 0.987 & 0.000 & 0.000 & 0.532 & 0.026 \\
XGB & smartwatch\_only & (4,3,7) & 216 & 0.005 & 0.176 & 0.000 & 0.987 & 0.000 & 0.000 & 0.247 & 0.010 \\
XGB & all & (4,3,14) & 427 & 0.005 & 0.148 & 0.000 & 0.961 & 0.000 & 0.000 & 0.528 & 0.016 \\
XGB & environment\_only & (4,3,14) & 171 & 0.005 & 0.125 & 0.000 & 0.987 & 0.000 & 0.000 & 0.371 & 0.012 \\
XGB & no\_environment & (4,3,14) & 315 & 0.005 & 0.093 & 0.000 & 0.842 & 0.000 & 0.000 & 0.323 & 0.011 \\
XGB & no\_peakflow & (4,3,14) & 355 & 0.005 & 0.060 & 0.000 & 0.984 & 0.000 & 0.000 & 0.595 & 0.022 \\
XGB & no\_smartinhaler & (4,3,14) & 400 & 0.005 & 0.116 & 0.000 & 0.969 & 0.000 & 0.000 & 0.599 & 0.022 \\
XGB & no\_smartwatch & (4,3,14) & 270 & 0.006 & 0.199 & 0.000 & 0.935 & 0.000 & 0.000 & 0.541 & 0.016 \\
XGB & peakflow\_only & (4,3,14) & 131 & 0.014 & 0.757 & 0.000 & 0.751 & 0.000 & 0.000 & 0.394 & 0.013 \\
XGB & smartinhaler\_only & (4,3,14) & 86 & 0.006 & 0.282 & 0.000 & 0.987 & 0.000 & 0.000 & 0.656 & 0.035 \\
XGB & smartwatch\_only & (4,3,14) & 216 & 0.005 & 0.176 & 0.000 & 0.987 & 0.000 & 0.000 & 0.408 & 0.013 \\
XGB & all & (4,7,0) & 427 & 0.010 & 0.510 & 0.000 & 0.920 & 0.000 & 0.000 & 0.487 & 0.016 \\
XGB & environment\_only & (4,7,0) & 171 & 0.006 & 0.173 & 0.000 & 0.977 & 0.000 & 0.000 & 0.297 & 0.011 \\
XGB & no\_environment & (4,7,0) & 315 & 0.007 & 0.297 & 0.000 & 0.865 & 0.000 & 0.000 & 0.392 & 0.012 \\
XGB & no\_peakflow & (4,7,0) & 355 & 0.007 & 0.327 & 0.000 & 0.990 & 0.667 & 0.400 & 0.712 & 0.250 \\
XGB & no\_smartinhaler & (4,7,0) & 400 & 0.009 & 0.475 & 0.000 & 0.863 & 0.000 & 0.000 & 0.449 & 0.013 \\
XGB & no\_smartwatch & (4,7,0) & 270 & 0.006 & 0.248 & 0.000 & 0.961 & 0.000 & 0.000 & 0.522 & 0.017 \\
XGB & peakflow\_only & (4,7,0) & 131 & 0.029 & 0.896 & 0.000 & 0.842 & 0.017 & 0.200 & 0.269 & 0.017 \\
XGB & smartinhaler\_only & (4,7,0) & 86 & 0.031 & 0.899 & 0.000 & 0.943 & 0.053 & 0.200 & 0.637 & 0.113 \\
XGB & smartwatch\_only & (4,7,0) & 216 & 0.009 & 0.450 & 0.000 & 0.982 & 0.000 & 0.000 & 0.518 & 0.047 \\
XGB & all & (4,7,7) & 427 & 0.011 & 0.574 & 0.000 & 0.909 & 0.000 & 0.000 & 0.449 & 0.015 \\
XGB & environment\_only & (4,7,7) & 171 & 0.006 & 0.173 & 0.000 & 0.981 & 0.000 & 0.000 & 0.211 & 0.010 \\
XGB & no\_environment & (4,7,7) & 315 & 0.007 & 0.267 & 0.000 & 0.855 & 0.000 & 0.000 & 0.338 & 0.012 \\
XGB & no\_peakflow & (4,7,7) & 355 & 0.008 & 0.416 & 0.000 & 0.984 & 0.333 & 0.200 & 0.702 & 0.293 \\
XGB & no\_smartinhaler & (4,7,7) & 400 & 0.007 & 0.332 & 0.000 & 0.906 & 0.000 & 0.000 & 0.422 & 0.014 \\
XGB & no\_smartwatch & (4,7,7) & 270 & 0.007 & 0.327 & 0.000 & 0.901 & 0.000 & 0.000 & 0.391 & 0.012 \\
XGB & peakflow\_only & (4,7,7) & 131 & 0.023 & 0.856 & 0.000 & 0.743 & 0.011 & 0.200 & 0.397 & 0.013 \\
XGB & smartinhaler\_only & (4,7,7) & 86 & 0.042 & 0.938 & 0.000 & 0.965 & 0.167 & 0.400 & 0.704 & 0.149 \\
XGB & smartwatch\_only & (4,7,7) & 216 & 0.008 & 0.371 & 0.000 & 0.981 & 0.000 & 0.000 & 0.548 & 0.033 \\
XGB & all & (4,7,14) & 427 & 0.005 & 0.089 & 0.000 & 0.895 & 0.029 & 0.200 & 0.476 & 0.022 \\
XGB & environment\_only & (4,7,14) & 171 & 0.006 & 0.218 & 0.000 & 0.981 & 0.000 & 0.000 & 0.279 & 0.011 \\
XGB & no\_environment & (4,7,14) & 315 & 0.006 & 0.134 & 0.000 & 0.906 & 0.000 & 0.000 & 0.347 & 0.012 \\
XGB & no\_peakflow & (4,7,14) & 355 & 0.005 & 0.089 & 0.000 & 0.975 & 0.250 & 0.400 & 0.780 & 0.327 \\
XGB & no\_smartinhaler & (4,7,14) & 400 & 0.005 & 0.059 & 0.000 & 0.834 & 0.018 & 0.200 & 0.429 & 0.016 \\
XGB & no\_smartwatch & (4,7,14) & 270 & 0.005 & 0.064 & 0.000 & 0.759 & 0.000 & 0.000 & 0.372 & 0.012 \\
XGB & peakflow\_only & (4,7,14) & 131 & 0.020 & 0.827 & 0.040 & 0.551 & 0.006 & 0.200 & 0.339 & 0.012 \\
XGB & smartinhaler\_only & (4,7,14) & 86 & 0.040 & 0.933 & 0.000 & 0.970 & 0.000 & 0.000 & 0.763 & 0.097 \\
XGB & smartwatch\_only & (4,7,14) & 216 & 0.007 & 0.307 & 0.000 & 0.967 & 0.111 & 0.200 & 0.597 & 0.094 \\
XGB & all & (4,14,0) & 427 & 0.034 & 0.847 & 0.043 & 0.943 & 0.056 & 0.250 & 0.472 & 0.025 \\
XGB & environment\_only & (4,14,0) & 171 & 0.008 & 0.333 & 0.000 & 0.934 & 0.048 & 0.250 & 0.520 & 0.035 \\
XGB & no\_environment & (4,14,0) & 315 & 0.005 & 0.000 & 0.000 & 0.943 & 0.000 & 0.000 & 0.634 & 0.019 \\
XGB & no\_peakflow & (4,14,0) & 355 & 0.019 & 0.721 & 0.037 & 0.951 & 0.067 & 0.250 & 0.590 & 0.074 \\
XGB & no\_smartinhaler & (4,14,0) & 400 & 0.028 & 0.809 & 0.040 & 0.934 & 0.048 & 0.250 & 0.563 & 0.030 \\
XGB & no\_smartwatch & (4,14,0) & 270 & 0.008 & 0.333 & 0.000 & 0.954 & 0.000 & 0.000 & 0.480 & 0.016 \\
XGB & peakflow\_only & (4,14,0) & 131 & 0.006 & 0.142 & 0.000 & 0.766 & 0.013 & 0.250 & 0.257 & 0.016 \\
XGB & smartinhaler\_only & (4,14,0) & 86 & 0.015 & 0.678 & 0.000 & 0.974 & 0.143 & 0.250 & 0.745 & 0.065 \\
XGB & smartwatch\_only & (4,14,0) & 216 & 0.005 & 0.000 & 0.000 & 0.917 & 0.000 & 0.000 & 0.777 & 0.046 \\
XGB & all & (4,14,7) & 427 & 0.020 & 0.727 & 0.000 & 0.959 & 0.083 & 0.250 & 0.807 & 0.089 \\
XGB & environment\_only & (4,14,7) & 171 & 0.012 & 0.546 & 0.000 & 0.941 & 0.056 & 0.250 & 0.448 & 0.030 \\
XGB & no\_environment & (4,14,7) & 315 & 0.005 & 0.000 & 0.000 & 0.950 & 0.000 & 0.000 & 0.793 & 0.050 \\
XGB & no\_peakflow & (4,14,7) & 355 & 0.020 & 0.732 & 0.000 & 0.962 & 0.091 & 0.250 & 0.772 & 0.278 \\
XGB & no\_smartinhaler & (4,14,7) & 400 & 0.020 & 0.732 & 0.000 & 0.947 & 0.062 & 0.250 & 0.710 & 0.042 \\
XGB & no\_smartwatch & (4,14,7) & 270 & 0.011 & 0.497 & 0.000 & 0.894 & 0.029 & 0.250 & 0.520 & 0.024 \\
XGB & peakflow\_only & (4,14,7) & 131 & 0.005 & 0.027 & 0.000 & 0.333 & 0.004 & 0.250 & 0.263 & 0.014 \\
XGB & smartinhaler\_only & (4,14,7) & 86 & 0.053 & 0.940 & 0.000 & 0.909 & 0.034 & 0.250 & 0.746 & 0.059 \\
XGB & smartwatch\_only & (4,14,7) & 216 & 0.005 & 0.000 & 0.000 & 0.917 & 0.000 & 0.000 & 0.793 & 0.051 \\
XGB & all & (4,14,14) & 427 & 0.013 & 0.579 & 0.000 & 0.914 & 0.071 & 0.500 & 0.717 & 0.091 \\
XGB & environment\_only & (4,14,14) & 171 & 0.008 & 0.295 & 0.000 & 0.898 & 0.032 & 0.250 & 0.488 & 0.025 \\
XGB & no\_environment & (4,14,14) & 315 & 0.005 & 0.000 & 0.000 & 0.889 & 0.056 & 0.500 & 0.764 & 0.065 \\
XGB & no\_peakflow & (4,14,14) & 355 & 0.008 & 0.350 & 0.000 & 0.905 & 0.065 & 0.500 & 0.818 & 0.358 \\
XGB & no\_smartinhaler & (4,14,14) & 400 & 0.008 & 0.339 & 0.000 & 0.942 & 0.105 & 0.500 & 0.715 & 0.111 \\
XGB & no\_smartwatch & (4,14,14) & 270 & 0.013 & 0.585 & 0.000 & 0.674 & 0.019 & 0.500 & 0.514 & 0.016 \\
XGB & peakflow\_only & (4,14,14) & 131 & 0.029 & 0.842 & 0.000 & 0.351 & 0.005 & 0.250 & 0.222 & 0.012 \\
XGB & smartinhaler\_only & (4,14,14) & 86 & 0.067 & 0.962 & 0.111 & 0.843 & 0.073 & 1.000 & 0.964 & 0.164 \\
XGB & smartwatch\_only & (4,14,14) & 216 & 0.006 & 0.027 & 0.000 & 0.911 & 0.069 & 0.500 & 0.924 & 0.538 \\
\end{longtable}


\paragraph{Leakage probe.}
Table~\ref{tab:v2_leakage} lists the per-configuration mean and
standard deviation of the validation ROC-AUC and PR-AUC across the
five shuffled-label refits.  The aggregate mean
($0.505 \pm 0.125$) is essentially at chance, indicating no
systematic information leakage from features that are functions of
labels or splits.

\begin{longtable}{lcccccc}
\caption{Leakage probe: 5 shuffled-label refits per configuration.  A leakage-free pipeline should show mean validation ROC-AUC near $0.5$.  Overall mean across all 81 configurations is $0.505$ $(\pm 0.125)$.}\label{tab:v2_leakage}\\
\toprule
Algo & $(T,L,W)$ & $n_{\text{shuf}}$ & Val ROC-AUC (mean$\pm$std) & Val PR-AUC (mean$\pm$std) \\
\midrule
\endfirsthead
\toprule
Algo & $(T,L,W)$ & $n_{\text{shuf}}$ & Val ROC-AUC (mean$\pm$std) & Val PR-AUC (mean$\pm$std) \\
\midrule
\endhead
\bottomrule
\endfoot
LR & (2,3,0) & 5 & 0.448 $\pm$ 0.257 & 0.010 $\pm$ 0.006 \\
LR & (2,3,7) & 5 & 0.394 $\pm$ 0.300 & 0.010 $\pm$ 0.007 \\
LR & (2,3,14) & 5 & 0.325 $\pm$ 0.138 & 0.007 $\pm$ 0.001 \\
LR & (2,7,0) & 5 & 0.399 $\pm$ 0.266 & 0.009 $\pm$ 0.004 \\
LR & (2,7,7) & 5 & 0.796 $\pm$ 0.122 & 0.066 $\pm$ 0.103 \\
LR & (2,7,14) & 5 & 0.761 $\pm$ 0.270 & 0.042 $\pm$ 0.035 \\
LR & (2,14,0) & 5 & 0.499 $\pm$ 0.255 & 0.014 $\pm$ 0.007 \\
LR & (2,14,7) & 5 & 0.564 $\pm$ 0.410 & 0.112 $\pm$ 0.217 \\
LR & (2,14,14) & 5 & 0.560 $\pm$ 0.304 & 0.020 $\pm$ 0.016 \\
LR & (3,3,0) & 5 & 0.400 $\pm$ 0.317 & 0.011 $\pm$ 0.007 \\
LR & (3,3,7) & 5 & 0.660 $\pm$ 0.300 & 0.024 $\pm$ 0.018 \\
LR & (3,3,14) & 5 & 0.586 $\pm$ 0.280 & 0.018 $\pm$ 0.013 \\
LR & (3,7,0) & 5 & 0.363 $\pm$ 0.338 & 0.012 $\pm$ 0.012 \\
LR & (3,7,7) & 5 & 0.228 $\pm$ 0.156 & 0.007 $\pm$ 0.001 \\
LR & (3,7,14) & 5 & 0.357 $\pm$ 0.254 & 0.009 $\pm$ 0.004 \\
LR & (3,14,0) & 5 & 0.481 $\pm$ 0.188 & 0.011 $\pm$ 0.004 \\
LR & (3,14,7) & 5 & 0.165 $\pm$ 0.152 & 0.007 $\pm$ 0.001 \\
LR & (3,14,14) & 5 & 0.471 $\pm$ 0.306 & 0.014 $\pm$ 0.008 \\
LR & (4,3,0) & 5 & 0.652 $\pm$ 0.169 & 0.018 $\pm$ 0.012 \\
LR & (4,3,7) & 5 & 0.771 $\pm$ 0.212 & 0.060 $\pm$ 0.068 \\
LR & (4,3,14) & 5 & 0.611 $\pm$ 0.318 & 0.049 $\pm$ 0.084 \\
LR & (4,7,0) & 5 & 0.486 $\pm$ 0.350 & 0.015 $\pm$ 0.010 \\
LR & (4,7,7) & 5 & 0.419 $\pm$ 0.459 & 0.050 $\pm$ 0.085 \\
LR & (4,7,14) & 5 & 0.484 $\pm$ 0.295 & 0.014 $\pm$ 0.011 \\
LR & (4,14,0) & 5 & 0.652 $\pm$ 0.168 & 0.019 $\pm$ 0.010 \\
LR & (4,14,7) & 5 & 0.539 $\pm$ 0.367 & 0.028 $\pm$ 0.032 \\
LR & (4,14,14) & 5 & 0.437 $\pm$ 0.207 & 0.011 $\pm$ 0.005 \\
RF & (2,3,0) & 5 & 0.552 $\pm$ 0.258 & 0.014 $\pm$ 0.010 \\
RF & (2,3,7) & 5 & 0.494 $\pm$ 0.368 & 0.022 $\pm$ 0.025 \\
RF & (2,3,14) & 5 & 0.561 $\pm$ 0.382 & 0.029 $\pm$ 0.033 \\
RF & (2,7,0) & 5 & 0.411 $\pm$ 0.419 & 0.018 $\pm$ 0.019 \\
RF & (2,7,7) & 5 & 0.397 $\pm$ 0.174 & 0.009 $\pm$ 0.003 \\
RF & (2,7,14) & 5 & 0.477 $\pm$ 0.169 & 0.011 $\pm$ 0.005 \\
RF & (2,14,0) & 5 & 0.645 $\pm$ 0.315 & 0.042 $\pm$ 0.049 \\
RF & (2,14,7) & 5 & 0.430 $\pm$ 0.190 & 0.010 $\pm$ 0.003 \\
RF & (2,14,14) & 5 & 0.495 $\pm$ 0.391 & 0.026 $\pm$ 0.026 \\
RF & (3,3,0) & 5 & 0.304 $\pm$ 0.150 & 0.007 $\pm$ 0.001 \\
RF & (3,3,7) & 5 & 0.393 $\pm$ 0.332 & 0.018 $\pm$ 0.027 \\
RF & (3,3,14) & 5 & 0.598 $\pm$ 0.439 & 0.039 $\pm$ 0.040 \\
RF & (3,7,0) & 5 & 0.511 $\pm$ 0.352 & 0.023 $\pm$ 0.028 \\
RF & (3,7,7) & 5 & 0.721 $\pm$ 0.233 & 0.072 $\pm$ 0.103 \\
RF & (3,7,14) & 5 & 0.645 $\pm$ 0.312 & 0.020 $\pm$ 0.010 \\
RF & (3,14,0) & 5 & 0.626 $\pm$ 0.344 & 0.215 $\pm$ 0.439 \\
RF & (3,14,7) & 5 & 0.326 $\pm$ 0.256 & 0.010 $\pm$ 0.006 \\
RF & (3,14,14) & 5 & 0.442 $\pm$ 0.292 & 0.015 $\pm$ 0.013 \\
RF & (4,3,0) & 5 & 0.607 $\pm$ 0.290 & 0.032 $\pm$ 0.045 \\
RF & (4,3,7) & 5 & 0.501 $\pm$ 0.415 & 0.030 $\pm$ 0.035 \\
RF & (4,3,14) & 5 & 0.515 $\pm$ 0.221 & 0.011 $\pm$ 0.004 \\
RF & (4,7,0) & 5 & 0.449 $\pm$ 0.361 & 0.015 $\pm$ 0.014 \\
RF & (4,7,7) & 5 & 0.400 $\pm$ 0.277 & 0.010 $\pm$ 0.007 \\
RF & (4,7,14) & 5 & 0.649 $\pm$ 0.298 & 0.052 $\pm$ 0.083 \\
RF & (4,14,0) & 5 & 0.609 $\pm$ 0.389 & 0.070 $\pm$ 0.104 \\
RF & (4,14,7) & 5 & 0.380 $\pm$ 0.241 & 0.010 $\pm$ 0.006 \\
RF & (4,14,14) & 5 & 0.358 $\pm$ 0.156 & 0.009 $\pm$ 0.002 \\
XGB & (2,3,0) & 5 & 0.447 $\pm$ 0.285 & 0.012 $\pm$ 0.009 \\
XGB & (2,3,7) & 5 & 0.433 $\pm$ 0.438 & 0.038 $\pm$ 0.060 \\
XGB & (2,3,14) & 5 & 0.516 $\pm$ 0.183 & 0.011 $\pm$ 0.004 \\
XGB & (2,7,0) & 5 & 0.508 $\pm$ 0.319 & 0.017 $\pm$ 0.016 \\
XGB & (2,7,7) & 5 & 0.526 $\pm$ 0.328 & 0.018 $\pm$ 0.016 \\
XGB & (2,7,14) & 5 & 0.522 $\pm$ 0.293 & 0.016 $\pm$ 0.013 \\
XGB & (2,14,0) & 5 & 0.679 $\pm$ 0.233 & 0.214 $\pm$ 0.440 \\
XGB & (2,14,7) & 5 & 0.621 $\pm$ 0.305 & 0.026 $\pm$ 0.020 \\
XGB & (2,14,14) & 5 & 0.425 $\pm$ 0.236 & 0.011 $\pm$ 0.005 \\
XGB & (3,3,0) & 5 & 0.385 $\pm$ 0.263 & 0.010 $\pm$ 0.009 \\
XGB & (3,3,7) & 5 & 0.493 $\pm$ 0.235 & 0.011 $\pm$ 0.005 \\
XGB & (3,3,14) & 5 & 0.473 $\pm$ 0.359 & 0.012 $\pm$ 0.007 \\
XGB & (3,7,0) & 5 & 0.454 $\pm$ 0.389 & 0.013 $\pm$ 0.008 \\
XGB & (3,7,7) & 5 & 0.775 $\pm$ 0.149 & 0.029 $\pm$ 0.015 \\
XGB & (3,7,14) & 5 & 0.580 $\pm$ 0.371 & 0.037 $\pm$ 0.045 \\
XGB & (3,14,0) & 5 & 0.666 $\pm$ 0.233 & 0.047 $\pm$ 0.068 \\
XGB & (3,14,7) & 5 & 0.373 $\pm$ 0.151 & 0.009 $\pm$ 0.002 \\
XGB & (3,14,14) & 5 & 0.616 $\pm$ 0.294 & 0.022 $\pm$ 0.016 \\
XGB & (4,3,0) & 5 & 0.459 $\pm$ 0.381 & 0.106 $\pm$ 0.220 \\
XGB & (4,3,7) & 5 & 0.577 $\pm$ 0.235 & 0.014 $\pm$ 0.009 \\
XGB & (4,3,14) & 5 & 0.605 $\pm$ 0.197 & 0.015 $\pm$ 0.008 \\
XGB & (4,7,0) & 5 & 0.468 $\pm$ 0.391 & 0.015 $\pm$ 0.012 \\
XGB & (4,7,7) & 5 & 0.396 $\pm$ 0.357 & 0.015 $\pm$ 0.018 \\
XGB & (4,7,14) & 5 & 0.598 $\pm$ 0.172 & 0.016 $\pm$ 0.013 \\
XGB & (4,14,0) & 5 & 0.533 $\pm$ 0.403 & 0.111 $\pm$ 0.218 \\
XGB & (4,14,7) & 5 & 0.410 $\pm$ 0.291 & 0.014 $\pm$ 0.012 \\
XGB & (4,14,14) & 5 & 0.325 $\pm$ 0.140 & 0.008 $\pm$ 0.002 \\
\end{longtable}


\paragraph{Comparison with v1.}  Only one $(T,L,W)$ cell overlaps
between the v1 questionnaire-event-episode experiment
(Section~\ref{sec:exp1}) and v2: $(3,7,7)$.  On that cell v1
reported PR-AUC $0.364$ and ROC-AUC $0.742$ on a $2$-way
$75/25$ patient split that allowed HPO on the same partition that
was later reported as the test set; v2 reports test
PR-AUC $\in [0.014, 0.023]$ and test
ROC-AUC $\in [0.46, 0.51]$ on the same cell for the three tabular
families (RF, XGB, LR) with the held-out $4$-patient test cohort.
We deliberately refrain from treating this as a fair head-to-head
comparison: the splits are different (no held-out validation in v1),
the test populations are different, and the test cohort carries on
the order of $2$--$5$ positives so the apparent v1 numbers were
resting on a patient-level effect that the new patient-disjoint test
split does not contain.  The purpose of v2 is methodological --- to
demonstrate that the v1 numbers, while internally consistent, are
inflated by HPO on the same partition that is later reported as the
test set --- not to claim numerical superiority on identical
conditions.  Full per-cell deltas are in
Table~\ref{tab:v2_vs_v1}.

\input{v2_tables/v1_vs_v2.tex}

\subsection{Reproducibility and audit trail}
\label{sec:v2-repro}

All v2 artefacts are emitted under \texttt{outputs/v2/tables/}:
\begin{itemize}
\itemsep0pt
\item \texttt{tune\_tabular\_v2\_\{trials,best,split\}.\{csv,json\}}
      --- HPO trials, winners, and the persisted three-way split.
\item \texttt{tune\_dl\_v2\_\{trials,best,multiseed,multiseed\_summary\}.csv}
      --- DL HPO and five-seed re-fit.
\item \texttt{sensor\_ablation\_v2.csv} --- nine-subset ablation
      for every $(T,L,W,\text{algo})$ tabular winner.
\item \texttt{leakage\_probe\_v2.\{csv,summary.csv\}} --- per-config
      shuffled-label re-fit and aggregate sanity check.
\item \texttt{final\_test\_v2\_\{tabular,dl,summary\}.csv} plus
      \texttt{final\_test\_v2\_predictions\_\{tabular,dl\}.csv}
      --- one-shot test metrics and per-sample probabilities.
\item \texttt{analysis\_v2\_\{tabular,dl\}\_ci.csv},
      \texttt{analysis\_v2\_v1\_vs\_v2.csv},
      \texttt{analysis\_v2\_sensor\_ablation\_topk.csv} ---
      offline post-hoc analysis (no training).
\end{itemize}
The entire pipeline is driven by
\texttt{notebooks/10\_event\_v2\_colab.ipynb} on Google Colab; a
fully executed copy is checked in as
\texttt{notebooks/10\_event\_v2\_colab\_fully\_ran.ipynb} alongside
the two result folders \texttt{v2\_results/} and
\texttt{v2\_results\_second\_run/} (first and second Colab runs
respectively).  All LaTeX tables in this section are regenerated by
running:
\begin{center}
\texttt{python scripts/v2/make\_v2\_latex\_tables.py}
\end{center}

\subsection{Limitations}
\label{sec:v2-limits}

Even with the corrected protocol, three caveats temper any clinical
reading of these numbers.  First, the test cohort contains only four
patients and two-to-five positive samples per cell; the bootstrap
$95\%$ CI on PR-AUC reaches the upper bound $1.0$ in every winning
configuration, so PR-AUC point estimates carry essentially no power.
ROC-AUC is more robust but still relies on a single patient-level
partition --- repeating the experiment under multiple seeds for the
split, rather than a single fixed seed, would tighten the confidence
intervals further.  Second, the strongest sensor-ablation result
depends on rescue-inhaler usage being recorded prospectively and
accurately; the dataset's smart-inhaler coverage is uneven across
patients, which limits external validity.  Third, the v2 task
definition emits positives only in the $[E-L,E-1]$ window: events
without an adequate observation history upstream (so the input
window cannot be constructed) are silently dropped, biasing the
cohort toward better-instrumented patients.  All of these limits
are made explicit in the per-cell sample counts
(Table~\ref{tab:v2_sample_counts}) so that any future re-evaluation
on a larger cohort can compare apples to apples.
